\documentclass[11pt]{article}

\usepackage[margin=1in]{geometry}
\usepackage{amsmath}
\usepackage{amssymb}
\usepackage{amsthm}
\usepackage{booktabs}
\usepackage{array}
\usepackage{longtable}
\usepackage[hidelinks]{hyperref}

\setlength{\emergencystretch}{3em}

\theoremstyle{plain}
\newtheorem{theorem}{Theorem}[section]
\newtheorem{lemma}[theorem]{Lemma}
\newtheorem{proposition}[theorem]{Proposition}
\newtheorem{corollary}[theorem]{Corollary}

\theoremstyle{definition}
\newtheorem{definition}[theorem]{Definition}
\newtheorem*{defn}{Definition}

\theoremstyle{remark}
\newtheorem{remark}[theorem]{Remark}
\newtheorem*{remk}{Remark}

\newcommand{\Z}{\mathbf{Z}}
\newcommand{\R}{\mathbf{R}}
\newcommand{\Cx}{\mathbf{C}}
\newcommand{\Gh}{\widehat{G}}
\newcommand{\tr}{\operatorname{tr}}
\newcommand{\Cled}{C_{\mathrm{led}}}
\newcommand{\Sgen}{S^{\mathrm{gen}}_{2}}
\newcommand{\Nd}{N_{d}}
\newcommand{\ph}{\widehat{\phi}}
\newcommand{\sinc}{\operatorname{sinc}}
\newcommand{\Marg}{\operatorname{Marg}}
\newcommand{\one}{\mathbf{1}}

\title{The two-moment certificate is robust under Rudnick--Sarnak-range cubic
augmentation with capacity control\\[2ex]
\large An exact absorption certificate for the two-moment system at bandwidth
one, with a hand-checkable witness}
\author{Kunal Tyagi\thanks{%
\texttt{hello.jay.tyagi@gmail.com}. \emph{Use of AI.} The mathematics in this
paper --- the derivations, the computations, the verification suite, the Lean
formalization, and the text --- was produced by Claude (Anthropic) working
under the author's direction; the author set the objectives, made the technical
and editorial decisions, and is responsible for the content. Claude is a tool
and is not an author. The complete working record is public:
\url{https://github.com/x67ai/riemann-rh-program}.}}
\date{August 26, 2026}

\begin{document}

\maketitle

\begin{abstract}
Alp\"oge and Furman prove unconditionally that more than two thirds of the
zeros of the Riemann zeta function are simple and lie on the critical
line~\cite{AF26}. Their argument rests on a two-moment certificate --- the
trace and the Frobenius norm of the Weil--Gabor compression, read against block
structure and integrality --- and that certificate is provably exhausted:
within the data class (tr, Frobenius, positive-index count, integer atoms) an
on-line double and a shallow off-line pair are indistinguishable, which pins
the distinct-zeros optimum at $\Nd/N = 5/6$ and the simple fraction at $2/3$.
A natural proposal to break this degeneracy is to adjoin, at a second bandwidth
$\lambda' = 1/2 + o(1)$ inside the Rudnick--Sarnak range, the unconditionally
computable signed third trace together with two spectral-capacity rows: an
all-$V$ eigenvalue-count ladder and a garnish-capacity fuzz row.

This paper proves, at the model level, and certifies computationally, that the
augmentation adds exactly nothing. The marginal value of the entire cubic block
over the two-moment baseline is $\delta_0 = 0$ identically, an exact LP
equality with an explicit, hand-checkable witness law. It holds over the full
adversary grid for the distinct-zeros benchmark ($940$ records,
$|\delta_0| < 10^{-12}$ everywhere, corner $5/6$); for the simple-fraction
benchmark (corner $2/3$) at its two decision cells together with cross-checks,
converged clean at the stopping criterion; and inside the near-CUE data class
licensed by the unconditional pair-correlation theorem of
Baluyot--Goldston--Suriajaya--Turnage-Butterbaugh~\cite{BGSTB24}, where the
absorption is by strict slack: there the cubic row does not bind.

The structural explanation is a set of exact identities proved here. On
isolated blocks the signed cubic row is the affine function $3F - 2$ of the
Frobenius row plus the second integrality level, hence blind to
doubles-versus-pairs at every depth. On the $\psi_1$-zero grid the
bandwidth-one Frobenius row degenerates exactly to the multiplicity count while
the half-bandwidth kernel stays position-sensitive, so the cubic budget is met
by arrangement at zero cost. Spectral escape is capped sharply at
$2\sqrt{2}\sqrt{\Cled \varepsilon_g}\,N = o(N)$, and a scalar cubic tail row
with any divergent cutoff is unconditionally vacuous.

The pair-interference channel, the one gap in the atom-only model theorem, is
closed by an integrality theorem throughout the admissible pair-depth family
($w \le 1$, plus deep probes at $w \in \{1.5, 2\}$): proved unconditionally for
single-pair columns to $w \le 2.8$, which covers that whole family, and for
multi-pair columns at $w \le 0.82$ ($w \le 0.98$ modulo one crowding cap whose
ledger constant is interval-certified; on the deepest sliver
$w \in (0.98, 1]$ the coverage is adversarial-numerical with the unconditional
$8/9$ backstop), with an unconditional $8/9$-strength spectral backstop at all
depths. Mixed-depth multi-pair columns involving the deep probes
$w \in \{1.5, 2\}$ rest on that backstop together with converged adversarial
pricing. The two-moment certificate, and with it the bandwidth-one ceiling, is
robust under Rudnick--Sarnak-range cubic augmentation with capacity control.
\end{abstract}

\section{Introduction}
\label{sec:intro}

\subsection{The two-moment certificate and its exhaustion}

Alp\"oge and Furman prove unconditionally~\cite{AF26} that more than two thirds
of the zeros of the Riemann zeta function are simple and on the critical line,
that $\Nd(T, 2T) \ge (5/6 - o(1)) N(T, 2T)$ zeros are distinct, and (Theorem~D,
window-optimized) that the simple-on-line proportion reaches
$2 - 1/c_1^{*} = 0.6725007\ldots$, by compressing the Weil explicit-formula form
to a critically sampled Gabor family and reading exactly two spectral moments of
the resulting Gram matrix $\Gh$: the trace $\tr \Gh = N(1 + o(1))$ and the
Frobenius moment $\tr \Gh^2 = (1/\lambda + \lambda/3) N (1 + o(1))$, both
evaluated unconditionally on the prime side for bandwidth $\lambda \le 1$. The
certificate is finite-dimensional linear algebra (the rank--trace inequality
``Lemma~R'') applied against the zero side's block structure: on-line points give
PSD rank-one atoms, off-line pairs give signature-$(1,1)$ hyperbolic blocks.

Two formalized results delimit this method exactly. First,
\texttt{lemmaR\_tight} (Lean,
\texttt{Zeta23/}\allowbreak\texttt{ZeroSide/}\allowbreak\texttt{TightMult.lean}
in the formalization accompanying~\cite{AF26}): on
configurations of integer-marked on-line atoms plus pair blocks, the two-moment
certificate is achieved with equality, and an on-line double (eigenvalue $+2$,
charge $4$) is indistinguishable from an off-line pair of depth $\to 0$
(eigenvalues $\pm 2$, charge $4$) --- no further inequality in
$(\tr, \text{Frobenius}, n_+, \text{integer atoms})$ can improve it. The
extremal configuration is $2/3$ simple zeros plus $1/6$ on-line doubles:
$\Nd/N = 5/6$, simple fraction $2/3$. Second, the bandwidth-one ceiling (Lean,
\texttt{Zeta23/PairCeiling/}): no certificate reading only bandwidth-one
form-factor data certifies a simple-zero proportion above $0.6818287$ (up to the
recorded stability terms), witnessed by an exact-rational $256$-periodic law; and
any trace/Frobenius certificate is capped below $0.8453 = 2 - 2/\sqrt{3}$ at any
Dirichlet-polynomial length, since $\kappa(\lambda) = 1/\lambda + \lambda/3 \ge
2/\sqrt{3}$ for all $\lambda > 0$.

\subsection{The cubic proposal}

Section~7.3 of~\cite{AF26} records the conditional template: under RH the third
trace $\tr \Gh^3$ is available (Hejhal~\cite{Hej94}, Rudnick--Sarnak~\cite{RS96})
and the certificate can
be run with the cubic weight
$\omega(m) = m/2 + (2m^2 - m^3)/18 + (4/9)\one_{m = 1}$, giving
$\Nd/N \ge 0.85082$ with the window $\cos(8s/5)$ (consuming also~\cite{BHB13}'s
$N^s/N \ge 19/27$ on RH). The proposal examined here is that a slice of this
becomes
unconditional: at bandwidth
$\lambda' = 1/2 + \vartheta \log\log T/\log T$, for a fixed offset exponent
$\vartheta > 0$ (strictly inside the $k = 3$
Rudnick--Sarnak range $X^3 \le T^{2 - \varepsilon}$), the signed cubic trace is
computable by the diagonal method of~\cite[\S5]{AF26}, and a Schatten-$3$
spectral-tail theorem would cap the spectral escape that the odd-moment no-go
of~\cite{AF26} (Section~7.5(e) of the 35-page version; the same item is
\S7.2(e) of v5, the two versions being distinguished in the References)
identifies, making the cubic row a usable sign-sensitive constraint.
That spectral-tail theorem has two forms, and this paper treats both: the
divergent-cutoff form in which the proposal stated it --- a single scalar tail
row $\sum_{|\lambda_i| \ge V_0} |\lambda_i|^3 = o(N)$ above one cutoff
$V_0 = V_0(T) \to \infty$ --- and the stronger all-$V$ eigenvalue-count ladder
$n(V) \le \Cled N V^{-4}$, holding at every threshold rather than above a single
cutoff ($\vartheta < 1$ strict). The same Chebyshev argument yields the ladder,
but only through a step --- from a bound above one cutoff to a bound at every
threshold --- that is not proved (Section~\ref{sec:flags}).
Theorems~\ref{thm:garnish} and~\ref{thm:capacity} treat the two forms in turn.
The proposal's separating premise was that a
hyperbolic pair block contributes $(+a)^3 + (-a)^3 = 0$ to the cubic while a
double contributes $+8$, so the two-moment-degenerate configurations would sit
at different corners of the augmented polytope.

Item~(e) of~\cite{AF26} is not silent on higher moments. In the 35-page version
it reads, verbatim:

\begin{quote}
``(e) The prime-side evaluation of $\tr \widetilde{G}^k$ by the diagonal method
of Section~5 (multiplicative relations among $k$ prime powers,
Montgomery--Vaughan for the rest) is available exactly in the Rudnick--Sarnak
range $k\lambda < 2$ [RS96]; for $\lambda \in (1/2, 1)$ this allows at most
$k = 3$ (and only for $\lambda < 2/3$), and an odd moment does not lower
$\Lambda_1(0)$. Thus, unconditionally, higher moments add nothing to the
$n_+$-bound on $(1/2, 1)$ (and Proposition~4.4 uses only the first two in any
case), while for $\lambda \le 1/2$, where many moments are available,
Proposition~7.4 makes them useless.''
\end{quote}

(In v5 the same item is compressed to ``\ldots at $X \asymp T$ this allows only
$k = 1$. Thus, unconditionally, higher moments add nothing.'') The remark is a
single-bandwidth statement, and it disposes of the single-bandwidth
question twice over, for two disjoint reasons: at bandwidth
$\lambda \in (1/2, 1)$ an odd moment does not move the Christoffel bound, and at
$\lambda \le 1/2$, where moments are plentiful, the certificate is worthless on
its own (its Proposition~7.4 cap is non-positive there). Alp\"oge and Furman
themselves label the remark informal.

The question addressed here is a two-bandwidth one, which that remark does not
reach: keep the bandwidth-one Frobenius baseline and adjoin a
second-bandwidth block at $\lambda' = 1/2$ --- the $\lambda'$-Frobenius budget,
the signed cubic equality, the all-$V$ count ladder, and the capacity fuzz row
(the widening of the cubic equality that Frobenius slack can purchase; rows
(R-6) and (R-5$'$) of Section~\ref{sec:rowsystem}).
Rows that are useless alone can bind when adjoined to a strong baseline, so what
the two-bandwidth question needs is a marginal-value computation rather than a
restatement; and the remark addresses neither the distinct-zeros benchmark, nor
the pair-block structure, nor the near-CUE data class, and exhibits no witness.
The answer obtained here is $\delta_0 = 0$ as an exact LP equality,
witness-backed, on both benchmarks, and by strict slack inside the near-CUE
class, together with the structural identities
(Sections~\ref{sec:structure}--\ref{sec:absorption}) and the capacity theorem
(Section~\ref{sec:capacity}) that explain it.

\paragraph{Related work.}
Gon\c{c}alves, de~Laat and Leijenhorst~\cite{GdLL25} run the
higher-level-correlation program against multiplicity directly, by semidefinite
programming over an auxiliary-function class: on GRH they obtain, from $3$-level
correlations, that at least $96.14\%$ of the zeros of a primitive $L$-function have
multiplicity at most $2$, and their $n = 3$ row pays where their $n = 2$ row does
not. That is consistent with the result proved here: theirs is a
single-bandwidth, conditional statement bounding \emph{multiplicity} at their
support parameter, whereas the one proved here is an unconditional statement
about the \emph{marginal value of adjoining a second bandwidth} to a fixed
bandwidth-one Frobenius baseline, at effective support $1$ (their $m = 2$ row),
measured against a fixed objective. The combinatorial half of their apparatus
and of this paper's is the same identity --- their equation~(5) is the Stirling
change of basis between the power sums $N_n = \sum_{\text{zeros}} m^n$ and the
falling-factorial counts $M_k = \sum_{\text{zeros}} m(m-1)\cdots(m-k+1)$, whose
$n = 3$ case is the atom half of Theorem~\ref{thm:blind} below --- so the atom
half is classical, and only the pair term $6\sum_p m_p^2 (m_p - 1) A_p^2$ is
new. The remark following their~(6) records that adapting their optimization to
bound $N_n$ directly did not improve on the bounds obtained through~(5); that
is a negative about problem \emph{formulation} rather than about the
information content of a higher correlation level. A closer structural
comparandum is sphere packing, where the three-point bound of Cohn, de~Laat and
Salmon~\cite{CdLS22} \emph{does} improve on the two-point Cohn--Elkies bound in
several dimensions: third-order augmentation buying exactly nothing here is
therefore a substantive negative, not a foregone one.

The marginal-value question is settled here by a linear program over an
adversary class --- marks up to a divergent alphabet $W$,
depth-parametrized shallow pairs, clustered sine-process null budgets, full
positional freedom, and an explicit garnish-capacity fuzz row derived from the
all-$V$ ladder. A strictly positive marginal value would be a gain to be
developed analytically; a marginal value of zero --- an \emph{absorption} of the
cubic block by the two-moment optimum --- sharpens the no-go. The value is zero.

\subsection{The marginal-value question and its answer}
\label{sec:marginalvalue}

Write $P_{\mathrm{base}}$ for the minimum of $E_w[\Nd]/N$ over laws $w$
(probability mixtures) of explicit $N$-periodic marked configurations satisfying
the two-moment data --- mass $N$, integrality, pair structure, bandwidth-one
Frobenius budget $(4/3) N (1 \mp \varepsilon)$ --- and $P_{\mathrm{full}}$ for
the same minimum after the entire $\lambda' = 1/2$ cubic block is adjoined (the
$\lambda'$-Frobenius budget $\kappa(1/2) N = (13/6) N$, the signed cubic
equality $c_3(1/2) N = 5N$, the all-$V$ count ladder, and the capacity fuzz
row). Both are linear programs, and Section~\ref{sec:model} fixes every term.
The question is whether the cubic block has strictly positive marginal value,
$\delta_0 := P_{\mathrm{full}} - P_{\mathrm{base}} > 0$.

Throughout, ``Rudnick--Sarnak-range cubic augmentation with capacity control''
(the title phrase) means exactly this
$\lambda' = 1/2 + o(1)$ system, with $\vartheta < 1$ strict --- the only point in
the RS range where both the signed cubic row and
the capacity ladder are unconditionally licensed; the behavior of the cubic row
at $\lambda' \ge 0.6$, inside the RS range but outside the proven ladder regime,
is the exploratory payoff curve of Section~\ref{sec:payoff}, not part of any
claim.

The answer, certified computationally over the decision grid and
explained by the model-level theorems below, is no: $\delta_0 = 0$
identically, not
merely within error bars, since the base-optimal law itself satisfies the entire
cubic block. The rest of the paper assembles the result: the structure theorems
that explain the absorption (Section~\ref{sec:structure}), the model-level corner
theorem with its pair-channel extension (Section~\ref{sec:absorption}), the
computational certification with the explicit witness
(Section~\ref{sec:certification}), the near-CUE result
(Section~\ref{sec:nearcue}), the simple-fraction benchmark
(Section~\ref{sec:p1}), the two capacity theorems that close the spectral-escape
routes (Section~\ref{sec:capacity}), and the scope map of what is exact,
what is a bounded residual, and what remains open (Section~\ref{sec:scope}).

\subsection{Headline results}
\label{sec:headline}

\begin{enumerate}
\item \textbf{Exact absorption with witness (Theorem~\ref{thm:corner} $+$
Section~\ref{sec:certification}).} Over the full decision grid ---
$940$ LP records spanning mark alphabets $W \in \{2, 3, 4, 6, 8\}$, ladder
constants $\Cled \in \{10, 100, 1000\}$, tolerances
$\varepsilon \in \{0.02, 0.05, 0.10\}$ (with the separate, labeled
tight-$\varepsilon$
runs extending to $\varepsilon = 0.002$), all fuzz modes, matched and asymptotic
budget centerings, $N = 64$ and $N = 128$ --- the marginal value of the cubic
block is $|\delta_0| < 10^{-12}$ at every record (maximum
$8.94 \times 10^{-13}$), and the joint optimum obeys the exact law
$P(\varepsilon) = 5/6 - (2/3)\varepsilon \to 5/6$. The absorbing law is
exhibited: three columns on the $65$-site $\psi_1$-zero grid, marks $\{1, 2\}$
only, no pairs, every row verified in exact rational arithmetic, the key entries
hand-checkable ($F_1 = 64$ exactly; $F' = 4128/33$).

\item \textbf{Both benchmarks.} The same absorption holds for the
simple-fraction objective: the cubic block moves the $p_1$ optimum by at most
$1.4 \times 10^{-12}$, converged at the stopping criterion (no improving
column in $S = 200$ pricing restarts, for the full and base systems alike), and
the
optimum is the exact analytic corner
$2 - (4/3)(1 + \varepsilon) \to 2/3$ (Section~\ref{sec:p1}).

\item \textbf{Near-CUE absorption by slack
(Section~\ref{sec:nearcue}).} Inside the pinned data class
$\bigl| E_w |c_j|^2 - j \bigr| \le 1$ (licensed unconditionally, with three
stated riders, by~\cite{BGSTB24}), the marginal value of the cubic block in that
pinned class is $\delta_0' = 0$, an exact LP equality at
$N = 64$ and $N = 128$, at both budget centerings --- and at the pinned optimum
only the pinning rows bind: the cubic row, the ladder, and even both Frobenius
rows are strictly interior.

\item \textbf{Cubic blindness (Theorem~\ref{thm:blind}).} On isolated blocks,
the signed cubic row equals $3F - 2$ per zero on every multiplicity-$1$ pair at
every depth and on every mark-$\{1, 2\}$ atom:
$C - 3F + 2\,\mathrm{Mass} = \sum m(m-1)(m-2) + 6 \sum m_p^2 (m_p - 1) A_p^2$
exactly. The ``$+8$ vs $0$'' separating premise of the proposal is
false: a pair's cubic charge is
$2m^3(1 + 3A^2) \ge 8m^3$ at every depth. All discriminating power of the cubic
row lives in marks $\ge 3$ and in clustering cross-terms --- and the cross-terms
are a free resource for the adversary (Theorem~\ref{thm:parseval}).

\item \textbf{Closed forms and the $m_2 = \kappa$ identity
(Theorem~\ref{thm:closedforms}).} The sine-process null budgets are
$m_2(\lambda) = 1/\lambda + \lambda/3$ and $m_3(\lambda) = 1 + 1/\lambda^2$ for
the flat window, all $0 < \lambda \le 1$; and $m_2(\lambda)$ coincides exactly
with the unconditional prime-side $\kappa(\lambda)$ of~\cite{AF26}: the sine
null's second-moment budget is the unconditional budget, on the whole band.

\item \textbf{Capacity cap (Theorem~\ref{thm:capacity}).} Under the all-$V$
ladder and a Frobenius slack $\varepsilon_g$, the maximal cubic charge
purchasable by spectral escape is exactly
$2\sqrt{2}\sqrt{\Cled \varepsilon_g}\,N = o(N)$ --- sharp, with the extremal
profile exhibited. The constant is confirmed by three independent derivations,
and it supersedes two earlier estimates, the two-term dyadic
$2\sqrt{C\varepsilon}$ and the point-mass profile's
$(5/\sqrt{3})\sqrt{C\varepsilon}$.

\item \textbf{Divergent-cutoff vacuity (Theorem~\ref{thm:garnish}).} A scalar
Schatten-$3$ tail row with any divergent cutoff is unconditionally absorbed by an
explicit ``garnish'' construction --- isolated on-line atoms placed below the
tail row's cutoff --- that lowers both objectives; this half of the no-go needs
no linear program at all.

\item \textbf{Pair-interference closure
(Sections~\ref{sec:fractional}--\ref{sec:pairchannel}).} The one channel outside
the atom-only sandwich --- off-line pairs pushing the bandwidth-one Frobenius
below the mark-accounting floor --- is closed for the admissible pair class
(Section~\ref{sec:adversary}) throughout its depth family: proved
unconditionally for all single-pair configurations to depth $w \le 2.8$
(covering every admissible depth, including both
deep probes, with $\ge 34\%$ margin), for all multi-pair configurations to
$w \le 0.82$ ($w \le 0.98$ modulo the crowding cap, ledger constant
interval-certified at $0.98465$; on $w \in (0.98, 1]$ the ledger fails ---
certified --- and coverage is adversarial numerics plus the $8/9$ backstop), and
for equal-depth pair-only configurations at every depth. Multi-pair
configurations at mixed depths involving the deep probes $w \in \{1.5, 2\}$ ---
admissible under Section~\ref{sec:adversary} but beyond the theorems' range ---
are covered by adversarial numerics plus the unconditional $8/9$ backstop
(see~(O2),
Section~\ref{sec:flags}). An explicit fractional-mark counterexample proves the
closure is intrinsically an integrality theorem.
\end{enumerate}

\subsection{What the cubic row reads, and the scope of the results}

The cubic row is a nonlinear spectral reading of bandwidth-one
pair-correlation data with quantitative high-spectrum capacity control. The
third trace at bandwidth $\lambda'$ carries total Fourier
support $(k - 1)\lambda' = 1 + o(1)$, not $3\lambda'$; its genuinely cubic prime
resonance is the absolute constant
$\sum_p (\log p)^3/(p - 1)^2 = 2.315762$ (evaluated by sieving to
$3 \times 10^6$), contributing $O(1)$ to $\tr \Gh^3$,
i.e.\ $O(1/N)$ per zero. No claim is made here that the cubic block evades the
Alternative Hypothesis, or that it consumes data beyond the bandwidth-one class
plus $o(1)$; for what the currently known band-limited $n$-level data provably
cannot do against AH, see Lagarias--Rodgers~\cite{LR20} and Tao's independent
treatment cited there. The no-go proved here is precisely that this
nonlinear reading, with the full capacity apparatus attached, adds nothing to
the linear reading at the proven operating point. Formally, a cubic row is
outside the PairCeiling certificate class as defined; the content of this paper
is that leaving the class buys nothing.

Every theorem in this paper is a statement about the model --- the explicit
configuration classes, kernels, and row systems of Section~\ref{sec:model} ---
not about the zeros of zeta. The bridge from the
model rows to unconditional statements about zeta is the sampling bookkeeping
of~\cite{AF26} plus the licensing riders of~\cite{BGSTB24}, stated in full in
Section~\ref{sec:license}. The exact zero $\delta_0 = 0$ is proved and
certified at the flat/flat window configuration and is specific to it: with the
Montgomery--Taylor cosine window at
$\lambda = 1$ the grid degeneracy breaks and the statement is a bounded
residual of order $6 \times 10^{-5}$ (Section~\ref{sec:scope}). What is
established is therefore not a universal ``$\delta_0 = 0$'': an exact zero on
the decision grid and in the near-CUE class; bounded residuals on the window
family, each below the $10^{-2}$ scale at which the question was posed; a proven
$o(N)$ cap on spectral escape everywhere; and an exploratory payoff curve
outside the proven ladder regime (Section~\ref{sec:payoff}).

\section{The model}
\label{sec:model}

This section fixes the model in which every theorem of Sections~3--7 lives: the
master trace formula, its finite-circle discretization, the adversary class, and
the row system.

\subsection{Window and kernel}

Fix a window profile $v \ge 0$ on $[-1/2, 1/2]$, even, with $\int v > 0$ (the
scale-free profile of the taper-squared $\phi^2$ of Alp\"oge and
Furman~\cite{AF26}; their formalized taper class is $C^3$). The \emph{normalized
sampling kernel} at bandwidth $\lambda > 0$ is the entire function
\[
  \psi_\lambda(x)
  = \frac{\displaystyle\int_{-1/2}^{1/2} v(\sigma)\,
      e^{2\pi i \lambda \sigma x}\, d\sigma}{\displaystyle\int v}\,,
\]
so $\psi_\lambda(0) = 1$, with Fourier transform $u$ supported in
$|\alpha| \le \lambda/2$ (Montgomery $\alpha$-units), $\int u = 1$. For the flat
window $v = 1$:
\[
  \psi_\lambda(x) = \sinc(\pi \lambda x)
  = \frac{\sin(\pi \lambda x)}{\pi \lambda x}\,,
  \qquad
  u = \frac{1}{\lambda}\,\one_{[-\lambda/2,\, \lambda/2]}\,.
\]
Positions $x$ are in mean-gap units (mean zero spacing $1$); complex positions
(off-line zeros) are allowed, $\psi_\lambda$ being entire.

\subsection{Continuum model: configurations and the master trace formula}
\label{sec:continuum}

A \emph{marked configuration} is a finite multiset of points $z$ with positions
$\gamma_z$ (complex allowed) and integer marks $m_z \ge 1$; off-line points occur
only in conjugate pairs $\{\gamma + i d, \gamma - i d\}$ with common depth and
common mark (functional-equation symmetry). The \emph{master trace formula}
defines the $k$-th trace at bandwidth $\lambda$:
\[
  \tr G^k := \sum_{z_1, \ldots, z_k}
    \Bigl( \prod_j m_{z_j} \Bigr)\,
    \psi_\lambda(\gamma_{z_1} - \gamma_{z_2})\,
    \psi_\lambda(\gamma_{z_2} - \gamma_{z_3}) \cdots
    \psi_\lambda(\gamma_{z_k} - \gamma_{z_1}),
\]
summed over ordered $k$-tuples of configuration points. This is the Poisson
sampling identity of~\cite{AF26} in normalized units (units in which an isolated
on-line zero of mark $m$ has eigenvalue exactly $m$), and it is realized by an
actual Hermitian matrix: Section~\ref{sec:pairblock} constructs the block
matrices explicitly where spectral (not just trace) data are needed. We write
$F = \tr G^2$ (Frobenius row) and $C = \tr G^3$ (signed cubic row); $F_1$ at
bandwidth $1$, $F'$ and $C'$ at bandwidth $\lambda'$. This single formula
generates every trace row on both sides --- null budgets and configuration
demands --- including all clustering cross-terms and all depth effects; no
isolated-block idealization enters the row system.

\subsection{Finite-circle model (the LP's discretization)}
\label{sec:finitecircle}

Period $N$ (zeros per period) on the circle of circumference $N$; configurations
are $N$-periodic, described by one period. For the flat window at bandwidth
$\lambda$ the harmonic set is $B = \{j \in \Z : |j| \le \lambda N/2\}$, with
$M = |B|$ harmonics and uniform weights $u_j = 1/M$. With the form factor
\[
  c_s = \sum_z m_z\, e^{-2\pi i s \theta_z/N}
\]
(a conjugate pair at $\theta \pm i d$ contributes
$2m \cosh(2\pi s d/N)\, e^{-2\pi i s \theta/N}$), the trace rows are
\begin{align*}
  \tr \Gh &= \sum_z m_z && \text{(mass)},\\
  \tr \Gh^2 &= \sum_s W_2(s)\, |c_s|^2,
    && W_2(s) = \sum_{j \in B,\; s - j \in B} u_j u_{s-j},\\
  \tr \Gh^3 &= \sum_{k_1, k_2} U_3(k_1, k_2)\, c_{k_1} c_{k_2} c_{-k_1-k_2},
    && U_3(k_1, k_2) = \sum_j u_{j+k_1} u_{j+k_1+k_2} u_j .
\end{align*}
Eigenvalue data (ladder counts, inertia) come from the Hermitian frequency matrix
$H[j, j'] = \sqrt{u_j u_{j'}}\, c_{j - j'}$ ($j, j' \in B$), whose $k$-th trace
equals the $k$-th trace row exactly (alias-free tight-frame compression). At the
primary geometry $N = 64$, $\lambda = 1$: $M = 65$; $\lambda' = 1/2$: $M' = 33$.
In position space, for a symmetric band,
$\tr \Gh^2 = \sum_{z, z'} m_z m_{z'} D(\theta_z - \theta_{z'})^2$ with
$D(x) = \sum_j u_j e^{-2\pi i j x/N}$ real, even, $D(0) = 1$ --- the discrete
analog of the continuum $\tr G^2 = \sum m\, m'\, \psi_\lambda^2$.

\subsection{Adversary class, laws, and the two benchmarks}
\label{sec:adversary}

A \emph{column}, also called an \emph{admissible} configuration below, is an
explicit $N$-periodic marked configuration: on-line atoms
(positions $\theta_i \in [0, N)$ --- positions are written $\theta$ throughout,
not to be confused with the offset exponent $\vartheta$ of
Section~\ref{sec:intro} --- integer marks $1 \le m_i \le W$) and off-line
pairs (position, dimensionless depth $w$ in the admissible pair-depth family
$(0, 1] \cup \{1.5, 2\}$, the last two values being the deep probes, pair
multiplicity $m$), with mass
$\sum m_i + 2 \sum m_p = N$ per period. A \emph{law} is a probability mixture
$w_c \ge 0$, $\sum_c w_c = 1$ over columns; every row is $E_w[\mathrm{row}(c)]$.
The two objectives:
\begin{align*}
  \Nd(c) &= \#\mathrm{atoms} + 2\#\mathrm{pairs}
    && \text{(distinct zeros; only on-line multiples reduce it)},\\
  p_1(c) &= \bigl( \#\text{mark-}1\ \mathrm{atoms}
      + 2\#\text{mult-}1\ \mathrm{pairs} \bigr)/N
    && \text{(simple fraction)}.
\end{align*}

\paragraph{Realizability.}
Imposed: explicit finite marked multisets (never bare correlation data); integer
marks; conjugate-pair structure; mass $N$; mixtures. \emph{Necessity:} an
adversary allowed to post aggregated moment data directly could claim values
realized by no point configuration (unconstrained cross-term variables can leave
the convex hull of realizable moment pairs), making an ``absorption'' vacuous;
this is the truncated-moment realizability problem of Kuna, Lebowitz and
Speer~\cite{KLS07,KLS11}, of which the finite-circle argument here is a special
case. Integrality, pair symmetry, and density are properties the true zero
multiset provably has, and dropping any of them hands the adversary provably
impossible worlds: with fractional marks even the $5/6$ baseline is false, since
mass $N$ of mark-$4/3$ atoms on the grid of Theorem~\ref{thm:parseval} has
$F_1 = (4/3) N$ budget-tight and $\Nd = (3/4) N < (5/6) N$
(Section~\ref{sec:fractional} exhibits the subtler pair-channel form of the same
failure). \emph{Sufficiency:} nothing further may be imposed, because every row of
the system is a theorem valid for arbitrary configurations with these properties
--- trace identities are linear algebra plus the sampling identity, the budgets
are the unconditional prime-side equalities, the ladder is the all-$V$
spectral-tail statement,
inertia is Sylvester; imposing anything more (e.g.\ pointwise GUE $3$-level
correlations) would smuggle in precisely the unproven correlation knowledge whose
absence the row system respects. Periodicity and the finite depth/mark grids are
discretizations, refined under a stability requirement
($N = 64 \to 128$, $W$-scan, depth-grid halving), not constraints.

\subsection{The row system}
\label{sec:rowsystem}

Baseline (two-moment / \texttt{lemmaR\_tight} data): (R-1) mass $N$ per period;
(R-2) the bandwidth-one
Frobenius budget $E_w[F_1] \in (4/3) N (1 \mp \varepsilon)$; (R-3) integrality
and pair structure, imposed by construction; (R-4) the objective. In the
matched-geometry variant the finite-$N$ CUE-sampled centers $B_1$, $B_2$, $B_3$
replace $(4/3)N$, $\kappa(1/2)N$ and $c_3(1/2)N$ respectively; their $N = 64$
values are given in Appendix~\ref{app:witness}. The cubic block at
$\lambda' = 1/2$ adds:
\begin{itemize}
\item[(R-5)] $\lambda'$-Frobenius:
  $E_w[F'] \in \kappa(1/2) N (1 \mp \varepsilon)$, $\kappa(1/2) = 13/6$;
\item[(R-6)] signed cubic with capacity coupling:
  $\bigl| E_w[C'] - c_3(1/2) N \bigr|
   \le \varepsilon_3 N + 2\sqrt{2}\sqrt{\Cled g_F}\, N$, $c_3(1/2) = 5$, the
  cubic tolerance being tied to the Frobenius one, $\varepsilon_3 = \varepsilon$.
  In the decoupled (parametric) fuzz mode the capacity term is replaced by a
  fixed cubic slack $\Gamma N$ with $\Gamma \in \{0, 0.05, 0.10, 0.25\}$, and the
  shadow price $d\delta_0/d\Gamma$ is reported;
\item[(R-5$'$)] garnish Frobenius budget $g_F \ge 0$ charged against the
  Frobenius slack;
\item[(R-7)] all-$V$ count ladder:
  $E_w[n(V)] \le \Cled N V^{-4}$ on the $V$-grid $\{2, 3, 4, 6, 8, 12, 16\}$;
\item[(R-8)] inertia: $n_-(c)$ is at most the number of pairs in $c$ (automatic
  on explicit spectra; asserted);
\item[(R-9)] Frobenius split, power-mean, H\"older-with-$o(N)$: automatically
  satisfied by explicit spectra (on
  explicit configurations the linear-algebra-true rows cannot be violated, so
  their entire content is enforced leak-free by construction).
\end{itemize}

The budget centers are the clustered sine-process null moments
(Section~\ref{sec:sine}) --- never isolated-block values --- because at
$\lambda' = 1/2$ the window width is two mean gaps and the null itself is
cross-term-dominated: cross-terms (occupancy $\ge 2$) carry $54\%$ of the
Frobenius budget and $80\%$ of the cubic budget (analytic fractions $7/6$ of
$13/6$ and $4$ of $5$; null MC $52.6\%\,/\,79.0\%$ at $n = 64$). The decision
rule: absorption iff
$\delta_0 = P_{\mathrm{full}} - P_{\mathrm{base}} = 0$ within LP tolerance at the
decision points with an explicit witness law; the witness must satisfy every row
at centered budgets with $E[\Nd]/N \le 5/6 + 0.01$. The objective, the adversary
class, the rows, the decision rule and the stopping criterion were fixed in
writing before the computation was run.

\section{Structure theorems}
\label{sec:structure}

On isolated blocks the signed cubic row is an affine function of the mass and
Frobenius rows plus the second integrality level, so it cannot see
doubles-versus-pairs at any depth (Theorem~\ref{thm:blind}). The null budgets
are then computed in closed form, yielding the identities that organize the
entire margin structure, and the grid-decoupling theorem shows that the
cross-term channel is available to the adversary at no cost. Every proof in this
section is complete; the numerical confirmations
(\texttt{results/a4-no-go/}\allowbreak\texttt{verify/verify\_t1.py}--\texttt{verify\_t4\_mc.py})
are summarized at the point of use.

\subsection{The exact pair block}
\label{sec:pairblock}

\paragraph{Setting.}
Windows $\tau_k = k(2\pi/L)$, $k \in \Z$ (critical spacing); taper $\phi$ real,
even, continuous, supported in $[-L/2, L/2]$, $\phi \in C^2$ (the class
formalized in~\cite{AF26} is $C^3$ --- more than enough). Fourier convention
$\ph(z) = \int \phi(s) e^{-i s z}\, ds$ (entire, by Paley--Wiener);
$\Phi(z) = \int \phi(s)^2 e^{-i s z}\, ds$; $a = L^{-1} \int \phi^2$, so
$\Phi(0) = L a$. Normalized units: matrices divided by $L^2 a$, in which an
isolated on-line atom of mark $m$ has eigenvalue $m$.

\begin{lemma}[Poisson sampling identity, complex arguments]
\label{lem:poisson}
For all complex $\tau, \tau'$:
\[
  \sum_{k \in \Z} \ph(\tau - \tau_k)\, \ph(\tau' - \tau_k)
  = L\, \Phi(\tau - \tau')\,,
\]
the sum converging absolutely and locally uniformly.
\end{lemma}

\begin{proof}
First real arguments. Fix $\tau, \tau' \in \R$ and set
$g(t) = \ph(\tau - t)\, \ph(\tau' - t)$; $\ph$ is real on $\R$ ($\phi$ real even)
and, since $\phi \in C^2$ with compact support,
$|\ph(x)| \le C/(1 + x^2)$, so $g \in L^1$ with
$|g(t)| \le C'/(1 + t^2)^2$ and the lattice sum converges absolutely. Each factor
$t \mapsto \ph(\tau - t)$ has Fourier transform
$e^{-i\tau\xi} 2\pi \phi(\xi)$ (inversion), supported in $[-L/2, L/2]$; hence
$\widehat{g} = (1/2\pi)(\text{convolution of the factor transforms})$ is
continuous, supported in $[-L, L]$, and vanishes at the endpoints $\pm L$ (the
convolution of two continuous functions supported in $[-L/2, L/2]$, evaluated at
$L$, integrates over the null set $\{L/2\}$). Poisson summation over the lattice
$(2\pi/L)\Z$ (classical form; $|g| + |\widehat{g}|$ summable on the respective
lattices):
\[
  \sum_k g(2\pi k/L)
  = \frac{L}{2\pi} \sum_{n \in \Z} \widehat{g}(nL)
  = \frac{L}{2\pi}\, \widehat{g}(0),
\]
all $n \ne 0$ terms vanishing by the support ($|nL| \ge L$, endpoint value $0$).
By Plancherel,
$\widehat{g}(0) = \int \ph(\tau - t) \ph(\tau' - t)\, dt
 = (1/2\pi) \int (2\pi\phi(\xi))^2 e^{-i(\tau - \tau')\xi}\, d\xi
 = 2\pi\, \Phi(\tau - \tau')$,
giving the identity on $\R^2$. For complex arguments:
$|\ph(x + iy)| \le e^{L|y|/2} C/(1 + x^2)$, so on compact subsets of $\Cx^2$ the
lattice sum converges absolutely and uniformly and is entire in $(\tau, \tau')$,
as is $L\Phi(\tau - \tau')$; the two entire functions agree on $\R^2$, hence
everywhere (identity theorem, one variable at a time).
\end{proof}

\begin{proposition}[exact off-line pair block --- the true Prop-4.1 spectrum]
\label{prop:pairblock}
Let the conjugate pair be $\{\gamma + i\delta, \gamma - i\delta\}$,
$\gamma \in \R$, depth $\delta > 0$, multiplicity $m$, with evaluation vector
$x = (x_k)$, $x_k = \ph(\gamma + i\delta - \tau_k)$, and block matrix
$M = m\,(x x^{\mathsf{T}} + \bar{x}\bar{x}^{\mathsf{T}})$ (transpose, not
conjugate-transpose: the two rank-one terms are the pair's two zeros). Then,
exactly:
\begin{enumerate}
\item[(i)] $x^{\mathsf{T}} x = L\,\Phi(0) = L^2 a$ --- real, independent of depth
  and position;
\item[(ii)] $|x|^2 = L\,\Phi(2 i \delta) = L^2 a\, A(\delta)$,
  $A(\delta) = \int \phi^2(s) \cosh(2\delta s)\, ds
   \big/ \int \phi^2(s)\, ds \ge 1$;
\item[(iii)] writing $x = p + i q$: $\langle p, q \rangle = 0$ exactly,
  $|p|^2 = L^2 a (1 + A)/2$, $|q|^2 = L^2 a (A - 1)/2$;
\item[(iv)] $M = 2m\,(p p^{\mathsf{T}} - q q^{\mathsf{T}})$, rank $\le 2$ real
  symmetric, with nonzero eigenvalues (normalized by $L^2 a$)
  \[
    m\,(1 + A(\delta)) \quad\text{and}\quad m\,(1 - A(\delta)),
  \]
  independent of the position $\gamma$;
\item[(v)] normalized charges: trace $2m$; Frobenius $2m^2(1 + A^2)$; signed
  cubic $2m^3(1 + 3A^2)$.
\end{enumerate}
In dimensionless depth $w := L\delta$, $A$ depends on the pair only through $w$,
with $A(0+) = 1$ and $A$ strictly increasing; for the flat taper
$A(w) = \sinh(w)/w$, with values $A(1/8) = 1.002606$, $A(1/4) = 1.010449$,
$A(1/2) = 1.042191$, $A(3/4) = 1.096422$, $A(1) = 1.175201$, $A(2) = 1.813430$.
\end{proposition}

\begin{proof}
(i) Lemma~\ref{lem:poisson} at $\tau = \tau' = \gamma + i\delta$:
$\sum_k \ph(\gamma + i\delta - \tau_k)^2 = L\Phi(0) = L \int \phi^2 = L^2 a$ ---
real, independent of $\gamma$ and $\delta$.
(ii) $\phi$ real and even gives $\ph(\bar{z}) = \overline{\ph(z)}$; hence
$|x_k|^2 = \ph(u - \tau_k)\, \ph(\bar{u} - \tau_k)$ with $u = \gamma + i\delta$,
and Lemma~\ref{lem:poisson} at $(u, \bar{u})$ gives
$|x|^2 = L\Phi(2i\delta) = L \int \phi^2(s)\cosh(2\delta s)\, ds$ ($\phi^2$
even); $A \ge 1$ by $\cosh \ge 1$, strictly increasing by strict monotonicity of
$\cosh$ on the support.
(iii) $x^{\mathsf{T}} x = |p|^2 - |q|^2 + 2i\langle p, q\rangle$; by (i) it is
real, so $\langle p, q \rangle = 0$ and $|p|^2 - |q|^2 = L^2 a$; with
$|p|^2 + |q|^2 = |x|^2 = L^2 a A$, solve.
(iv) $x x^{\mathsf{T}} + \bar{x}\bar{x}^{\mathsf{T}}
 = (p + iq)(p + iq)^{\mathsf{T}} + (p - iq)(p - iq)^{\mathsf{T}}
 = 2(p p^{\mathsf{T}} - q q^{\mathsf{T}})$; with $p \perp q$ the eigenvectors are
$p, q$ with eigenvalues $2m|p|^2 = m L^2 a (1 + A)$ and
$-2m|q|^2 = m L^2 a (1 - A)$; all others $0$; normalize. Position-independence:
the eigenvalues depend only on $|p|^2, |q|^2$, which depend only on $\delta$.
(v) Trace $m(1 + A) + m(1 - A) = 2m$; Frobenius
$m^2(1 + A)^2 + m^2(1 - A)^2 = 2m^2(1 + A^2)$; cubic
$m^3(1 + A)^3 + m^3(1 - A)^3 = 2m^3(1 + 3A^2)$. Flat-taper $A$:
$A(w) = \int_{-1/2}^{1/2} \cosh(2w\sigma)\, d\sigma = \sinh(w)/w$; the listed
values are its evaluations (they reproduce the tabulated values used in the
computation to their displayed precision, $\le 4 \times 10^{-6}$).
\end{proof}

\paragraph{Consistency and inertia.}
The two-point configuration $\{\gamma \pm i\delta\}$ under the master formula has
$\tr G^k = m^k \tr K^k$ with the $2 \times 2$ kernel
$K = \bigl[\begin{smallmatrix} 1 & A \\ A & 1 \end{smallmatrix}\bigr]$,
eigenvalues $1 \pm A$ --- the same charges. The matrix route (iv) additionally
certifies the inertia (one positive, one negative eigenvalue for $\delta > 0$)
consumed by the $n_-$ and ladder rows.

\paragraph{Block continuity, exact form.}
As $\delta \to 0+$, the spectrum tends to $(2m, -0)$, Frobenius to $2m^2$, cubic
to $8m^3$: a shallow pair is row-indistinguishable from an on-line double in
every row, continuously. Moreover the cubic charge $2m^3(1 + 3A^2) \ge 8m^3$ at
every depth ($A \ge 1$), increasing with depth: a deep pair never loses cubic
charge. The ``$+8$ vs $0$'' separation premise is false under the true block
law: it is the cubic bookkeeping that fails, the block spectrum
$m L^2 a (1 \pm A)$ of Proposition~\ref{prop:pairblock} being unaffected.

Numerical confirmation (\texttt{verify\_t3.py}): evaluation-vector
route with a $\cos$ taper, $L = 40$, $120001$ windows:
$|x^{\mathsf{T}} x / L^2 a - 1| \le 6.7 \times 10^{-16}$,
$|\langle p, q\rangle|/L^2 a \le 1.7 \times 10^{-14}$, block eigenvalues vs
$m(1 \pm A)$ to $7.8 \times 10^{-14}$, position-independent at machine precision
at every depth $w \in \{1/64, 1/8, 1/4, 1/2, 3/4, 1, 2\}$; shallow-pair cubic
charge $8.0005$ at $w = 1/64$ and $8.0313$ at $w = 1/8$, matching the
independently computed $8.0005 / 8.033$ --- never $0$ at any depth.

\subsection{The affine-plane identity and cubic blindness}

\begin{defn}[per-zero charges]
For an isolated block of total mass $\mu$ ($\mu = m$ for an atom of mark $m$;
$\mu = 2m$ for a pair of multiplicity $m$), let
$F = (\text{block Frobenius})/\mu$ and $c = (\text{block cubic})/\mu$. By
Proposition~\ref{prop:pairblock}(v) and the atom spectrum:
\begin{align*}
  \text{atom, mark } m: &\quad F = m, && c = m^2;\\
  \text{pair, mult } m,\ \text{depth}: &\quad F = m(1 + A^2),
    && c = m^2(1 + 3A^2).
\end{align*}
\end{defn}

\begin{proposition}[affine-plane identity]
\label{prop:affine}
\begin{align*}
  \text{atom of mark } m: &\quad c = 3F - 2 + (m-1)(m-2)\,;\\
  \text{pair of mult } 1,\ \text{any depth}: &\quad c = 3F - 2 \quad\text{exactly}\,;\\
  \text{pair of mult } m,\ \text{any depth}: &\quad
    c = 3F - 2 + (m-1)(m-2) + 3m(m-1)A^2 \,.
\end{align*}
In particular every mark-$\{1,2\}$ atom and every multiplicity-$1$ pair --- at
every depth --- lies on the same affine plane $c = 3F - 2$ in the per-zero
$(F, c)$ plane.
\end{proposition}

\paragraph{The atom half in the literature.}
For atoms the identity $c - 3F + 2 = (m-1)(m-2)$, and its summed form in
Theorem~\ref{thm:blind}, is the third-order Stirling change of basis between
power sums and falling factorials: $m(m-1)(m-2) = m^3 - 3m^2 + 2m$, the
coefficients being the signed Stirling numbers of the first kind $s(3,3) = 1$,
$s(3,2) = -3$, $s(3,1) = 2$. The inverse conversion,
$\sum m^n = \sum_k S(n,k) M_k$ with Stirling numbers of the second kind, is
equation~(5) of~\cite{GdLL25}, where it plays the same role --- including the
same consequence, that the falling-factorial term vanishes identically once
every multiplicity is at most $n-1$. The atom half is thus classical; what
Theorem~\ref{thm:blind} adds is the pair term
$6\sum_{\text{pairs}} m_p^2 (m_p - 1) A_p^2$, a statement about the finite-circle
block law of Proposition~\ref{prop:pairblock} with no counterpart in the
multiplicity-statistics literature.

\begin{proof}
Atom: $3F - 2 + (m-1)(m-2) = 3m - 2 + m^2 - 3m + 2 = m^2 = c$. Pair:
$c - 3F + 2 = m^2(1 + 3A^2) - 3m(1 + A^2) + 2
 = (m^2 - 3m + 2) + 3A^2(m^2 - m) = (m-1)(m-2) + 3m(m-1)A^2$. At $m = 1$ both
terms vanish identically in $A$ (every depth); at $m = 2$ the atom identity also
reads $c = 3F - 2$ (doubles are on the plane).
\end{proof}

\begin{theorem}[cubic blindness on isolated blocks]
\label{thm:blind}
Let the configuration be a disjoint union of isolated blocks (atoms of marks
$m_i$; pairs of multiplicities $m_p$ at depths with $A_p = A(\delta_p)$), with
total mass $\mathrm{Mass}$, Frobenius $F$, signed cubic $C$ (sums of block
charges). Then, exactly:
\[
  C - 3F + 2\,\mathrm{Mass}
  = \sum_{\mathrm{zeros}} m(m-1)(m-2)
    + 6 \sum_{\mathrm{pairs}} m_p^2 (m_p - 1) A_p^2 \,,
\]
where the first sum runs over the zero multiset (an atom of mark $m$ contributes
$m(m-1)(m-2)$ once; a pair of mult $m$ contributes twice that --- its two zeros).
Consequences:
\begin{enumerate}
\item On the class $\{\text{mark-}\{1,2\}\ \text{atoms}\} \cup
  \{\text{multiplicity-}1\ \text{pairs of arbitrary depth}\}$ --- the entire
  doubles-vs-pairs comparison class of \texttt{lemmaR\_tight} --- the right side
  vanishes identically: $C = 3F - 2\,\mathrm{Mass}$ identically on the
  class. The signed cubic row restricted to isolated blocks of this class is an
  exact affine combination of the Frobenius and mass rows: it separates nothing
  those two rows do not. An on-line double and an off-line pair of any depth are
  cubic-indistinguishable.
\item All discriminating power of the cubic row beyond the two-moment data
  therefore lives in exactly two places: (i) marks $\ge 3$ (and pair
  multiplicities $\ge 2$), through the nonnegative integrality terms on the
  right; (ii) violations of block isolation --- the clustering cross-terms of the
  master trace formula, which the identity does not constrain, and which
  Theorem~\ref{thm:parseval} shows to be adversary-tunable at zero cost.
\item The premise ``a hyperbolic pair block contributes $(+a)^3 + (-a)^3 = 0$ to
  $\tr G^3$ while a double contributes $+8$'' is false under the true block law:
  the pair's signed cubic is $2m^3(1 + 3A^2) \ge 8m^3$ at every depth. The
  sign-split ``different corners of the constraint polytope'' picture is
  superseded.
\end{enumerate}
\end{theorem}

\begin{proof}
Sum the per-block identities. An atom of mark $m$ has block charges
(mass $m$, Frobenius $m^2$, cubic $m^3$), contributing
$m^3 - 3m^2 + 2m = m(m-1)(m-2)$ to $C - 3F + 2\,\mathrm{Mass}$. A pair of mult
$m$ has block charges (mass $2m$, Frobenius $2m^2(1 + A^2)$, cubic
$2m^3(1 + 3A^2)$), contributing
\[
  2m^3(1 + 3A^2) - 6m^2(1 + A^2) + 4m
  = 2m(m-1)(m-2) + 6m^2(m-1)A^2 \,,
\]
exactly its two zeros' worth of the integrality sum plus the
multiplicity-weighted depth term. On marks $\{1,2\}$ and mult-$1$ pairs every
term vanishes. Consequence 3 is Proposition~\ref{prop:pairblock} ($A \ge 1$).
\end{proof}

In classical language this is the Chebyshev--Markov--Stieltjes phenomenon that an
odd moment does not improve a one-sided bound in a Chebyshev system (Karlin and
Studden~\cite{KS66}); Theorem~\ref{thm:blind} is its exact form for the block
class at issue here.

\begin{remk}[why the row system is clustered]
By consequence~1, a linear program built on isolated blocks over the
doubles-vs-pairs class returns a tautological answer along the plane
$c = 3F - 2$. The clustered null budgets and positional freedom of
Section~\ref{sec:rowsystem} are forced rather than conservative, and the
clustering cross-terms, the only live channel, are a free resource for the
adversary rather than a signal for the certificate.
\end{remk}

Numerical confirmation (\texttt{verify\_t3.py}): charges and the per-zero
identity for $m \in \{1, 2, 3, 5\}$ at all depths to $\le 9.1 \times 10^{-13}$;
atom identity exhaustive $m = 1, \ldots, 29$, exact; master-formula
$2 \times 2$ kernel route reproduces all three charges.

\subsection{Sine-process closed forms and the doubles-plane diagnostic}
\label{sec:sine}

The null budgets are the per-zero moments of the flat-window Gram over the
determinantal sine process --- the process with kernel
$S(x) = \sin(\pi x)/(\pi x)$, intensity $1$, and correlation functions (taken as
the null's definition)
\begin{align*}
  \rho_2(x_1, x_2) &= 1 - S(x_{12})^2,\\
  \rho_3(x_1, x_2, x_3) &= 1 - S(x_{12})^2 - S(x_{13})^2 - S(x_{23})^2
    + 2 S(x_{12}) S(x_{13}) S(x_{23}).
\end{align*}
The per-zero Gram moments at bandwidth $\lambda$ (flat window,
$\psi = \psi_\lambda$, $u = \widehat{\psi}$):
\begin{align*}
  m_2(\lambda) &:= 1 + \int \rho_2(0, x)\, \psi(x)^2\, dx \,,\\
  m_3(\lambda) &:= 1 + 3 \int \rho_2(0, x)\, \psi(x)^2\, dx + T_3(\lambda) \,,\\
  T_3(\lambda) &:= \iint \rho_3(0, x, x + y)\, \psi(x) \psi(y) \psi(x + y)\,
    dx\, dy \,,
\end{align*}
the $N \to \infty$ limits of $E[\tr G^k]/N$: the $k$-tuple sum of
Section~\ref{sec:continuum} splits over coincidence patterns --- all-equal (value
$1$), exactly-two-equal ($3$ ordered patterns at $k = 3$, each $\psi(x)^2$
against $\rho_2$), all-distinct ($\rho_3$ against the cyclic kernel product,
reparametrized by separations using the evenness of $\psi$). The sine process is
simple, so all marks are $1$.

\begin{theorem}[flat-window closed forms]
\label{thm:closedforms}
For the flat window and every $0 < \lambda \le 1$:
\[
  m_2(\lambda) = \frac{1}{\lambda} + \frac{\lambda}{3} \,,
  \qquad
  m_3(\lambda) = 1 + \frac{1}{\lambda^2} \,.
\]
\end{theorem}

\begin{proof}
\emph{Step 1 (Fourier-side reduction).} With $t(\alpha) = (1 - |\alpha|)_+$ (the
transform of $S^2$), $b = \one_{[-1/2, 1/2]}$ ($= \widehat{S}$), and the
convention $\widehat{f}(\alpha) = \int f(x) e^{-2\pi i \alpha x}\, dx$:
\begin{enumerate}
\item[(i)] $\int \psi^2\, dx = \int u^2\, d\alpha$ \hfill [Plancherel];
\item[(ii)] $\int S^2 \psi^2\, dx = \int t\, (u * u)\, d\alpha$
  \hfill [Plancherel on the product pair];
\item[(iii)] $\iint \psi(x) \psi(y) \psi(x + y)\, dx\, dy = \int u^3\, d\alpha$;
\item[(iv)] $\iint S(x)^2 \psi(x) \psi(y) \psi(x + y)\, dx\, dy
  = \int (t * u)\, u^2\, d\alpha$, and the same value for the $S(y)^2$ and
  $S(x+y)^2$ variants;
\item[(v)] $\iint S(x) S(y) S(x + y) \psi(x) \psi(y) \psi(x + y)\, dx\, dy
  = \int (b * u)^3\, d\alpha$.
\end{enumerate}
For (iii), two Plancherel steps with no delta calculus: with $y$ fixed,
$g(y) := \int \psi(x) \psi(x + y)\, dx$ is the Fourier transform of $u^2$ at $y$;
then $\int \psi(y) g(y)\, dy = \int \widehat{\psi}\, u^2 = \int u^3$ by the
multiplication identity $\int f \widehat{g} = \int \widehat{f} g$; Fubini is
licensed by
$|\psi(x)\psi(y)\psi(x+y)| \le C/((1+|x|)(1+|y|)(1+|x+y|))$, integrable on
$\R^2$. For (iv), the same two-step Plancherel with first factor
$S(x)^2 \psi(x)$ gives $\int (S^2\psi)^{\widehat{\ }}\, u^2 = \int (t * u) u^2$;
the $S(y)^2$ and $S(x+y)^2$ variants reduce to the same value under the
substitutions $(x, y) \mapsto (y, x)$ and $(x, y) \mapsto (x, -x-y)$ (Jacobian
$1$, integrand invariant). For (v), set $h = S\psi$, $\widehat{h} = b * u$, and
apply (iii) to $h$. Hence
\begin{align*}
  m_2 &= 1 + \int u^2 - \int t\,(u * u), &
  m_3 &= 1 + 3\Bigl[ \int u^2 - \int t\,(u * u) \Bigr] + T_3,\\
  T_3 &= \int u^3 - 3 \int (t * u) u^2 + 2 \int (b * u)^3. &&
\end{align*}

\emph{Step 2 (evaluation at $u = (1/\lambda)\one_{[-\lambda/2, \lambda/2]}$,
using $\lambda \le 1$ exactly where indicated).}
\begin{enumerate}
\item[(1)] $\int u^2 = 1/\lambda$.
\item[(2)] $(u * u)(\alpha) = (1/\lambda^2)(\lambda - |\alpha|)_+$; since
  $\lambda \le 1$, $t = 1 - |\alpha|$ on this support, so
  $\int t\,(u * u) = (2/\lambda^2) \int_0^\lambda (1 - a)(\lambda - a)\, da
   = (2/\lambda^2)[\lambda^2/2 - \lambda^3/6] = 1 - \lambda/3$. Hence
  $m_2 = 1 + 1/\lambda - (1 - \lambda/3) = 1/\lambda + \lambda/3$.
\item[(3)] $\int u^3 = 1/\lambda^2$.
\item[(4)] $\int (t * u) u^2
  = (1/\lambda^3) \iint_{|\alpha| \le \lambda/2,\ |\alpha - \beta| \le \lambda/2}
    t(\beta)\, d\beta\, d\alpha
  = (1/\lambda^3) \int t(\beta)(\lambda - |\beta|)_+\, d\beta
  = (1/\lambda^3)[\lambda^2 - \lambda^3/3] = 1/\lambda - 1/3$
  (reusing the integral of (2); valid for $\lambda \le 1$).
\item[(5)] $(b * u)(\alpha)$ is the trapezoid equal to $1$ on
  $|\alpha| \le (1 - \lambda)/2$, decaying linearly to $0$ at
  $|\alpha| = (1 + \lambda)/2$; hence
  $\int (b * u)^3 = (1 - \lambda) + 2\lambda \int_0^1 s^3\, ds
   = 1 - \lambda/2$.
\end{enumerate}
Assembling: $T_3 = 1/\lambda^2 - 3(1/\lambda - 1/3) + 2(1 - \lambda/2)
= 1/\lambda^2 - 3/\lambda + 3 - \lambda$, and
$m_3 = 1 + 3[1/\lambda - 1 + \lambda/3] + T_3 = 1 + 1/\lambda^2$.
\end{proof}

\begin{corollary}[anchor values, the doubles-plane diagnostic, the moment margin]
\label{cor:anchors}
\[
  m_2(1) = 4/3, \quad m_3(1) = 2, \quad m_2(1/2) = 13/6, \quad m_3(1/2) = 5.
\]
Define $G(\lambda) := m_3 - (3 m_2 - 2)$ (the null budget's excess over the
isolated-block doubles/pairs plane of Theorem~\ref{thm:blind}) and
$\Marg(\lambda) := 2 m_2 - m_3$ (the certificate margin
of~\cite[\S7.3]{AF26}). Then, flat window, on $(0, 1]$:
\begin{align*}
  G(\lambda) &= 3 + \frac{1}{\lambda^2} - \frac{3}{\lambda} - \lambda \,;
    & G(1) &= 0, & G(1/2) &= 1/2,\\
  \Marg(\lambda) &= \frac{2}{\lambda} + \frac{2\lambda}{3} - 1
    - \frac{1}{\lambda^2} \,;
    & \Marg(1) &= 2/3, & \Marg(1/2) &= -2/3,
\end{align*}
and $\Marg$ has a unique sign change on $(0, 1]$ (the interval on which the
closed forms are proved), at $\lambda_0 = 0.610511\ldots$ Each value is
load-bearing for the no-go:
\begin{enumerate}
\item \textbf{$G(1) = 0$ is the ``$2 = 2$'' saturation.} At bandwidth one the
  sine budget $m_3(1) = 2$ exactly equals the isolated-block cubic demand of the
  $5/6$-extremal ($2/3$ simple $+\ 1/6$ doubles: $(2/3)(1) + (1/3)(4)$
  mass-weighted $= 2$, on the plane of Theorem~\ref{thm:blind}). The bandwidth-one
  cubic row alone cannot cut the doubles corner --- the budget is exactly met.
\item \textbf{$G(1/2) = +1/2 > 0$.} At the operating bandwidth the null budget
  sits strictly above the doubles plane: a doubles-saturated adversary must
  supply $+1/2$ per zero of cubic through positional cross-terms --- which
  Theorem~\ref{thm:parseval}'s grid freedom supplies at zero cost (the
  witness of Section~\ref{sec:certification} does exactly this, sitting at the
  lower cubic budget edge).
\item \textbf{$\Marg(1/2) = -2/3 < 0$.} The $\omega(m)$-moment route's margin is
  negative at the proven operating point: a strictly positive marginal value
  could never have come from the mechanism of~\cite[\S7.3]{AF26} there, only
  from joint positional pricing. The sign change at $\lambda_0 = 0.6105$ is
  precisely where the exploratory scan saw $\delta_0$ turn positive (between
  $\lambda' = 0.55$ and $0.60$; Section~\ref{sec:payoff}), outside the proven
  ladder regime.
\end{enumerate}
\end{corollary}

\begin{proof}
Substitute Theorem~\ref{thm:closedforms}; the displayed formulas are algebraic
identities on $(0, 1]$. $G(1) = 3 + 1 - 3 - 1 = 0$;
$G(1/2) = 3 + 4 - 6 - 1/2 = 1/2$; $\Marg(1) = 2 + 2/3 - 1 - 1 = 2/3$;
$\Marg(1/2) = 4 + 1/3 - 1 - 4 = -2/3$. Uniqueness of the root:
$q(\lambda) = \lambda^2 \Marg(\lambda)
= 2\lambda + (2/3)\lambda^3 - \lambda^2 - 1$ has
$q' = 2 + 2\lambda^2 - 2\lambda > 0$ (negative discriminant), so $q$ is strictly
increasing with $q(0.55) < 0 < q(0.66)$; numerically the root is $0.610511$
(\texttt{brentq}, re-verified by \texttt{sympy} \texttt{nsolve}). The extremal
demand in item~1: per-zero $(F, c) = (1, 1)$ on simples and $(2, 4)$ on doubles;
mass-weighted cubic $(2/3)(1) + (1/3)(4) = 2 = m_3(1)$ and Frobenius
$(2/3)(1) + (1/3)(2) = 4/3 = m_2(1)$ --- both budgets exactly saturated at
$\lambda = 1$.
\end{proof}

\begin{corollary}[the $m_2 = \kappa$ identity: sine null $=$ unconditional budget]
\label{cor:m2kappa}
The closed form $m_2(\lambda) = 1/\lambda + \lambda/3$ coincides, at every
$\lambda \le 1$, with the value obtained by integrating the unconditional
Montgomery pair-correlation form factor of~\cite{BGSTB24},
$F(\alpha) = \delta_{\mathrm{D}}(\alpha) + |\alpha|$ (on $|\alpha| \le 1$),
$\delta_{\mathrm{D}}$ being the Dirac mass at the origin --- not the marginal
value $\delta_0$ --- against the flat-window pair kernel:
\[
  \frac{1}{\lambda^2} \int_{-\lambda}^{\lambda} (\lambda - |\alpha|)
    \bigl[ \delta_{\mathrm{D}}(\alpha) + |\alpha| \bigr]\, d\alpha
  = \frac{\lambda + \lambda^3/3}{\lambda^2}
  = \frac{1}{\lambda} + \frac{\lambda}{3}
  = m_2(\lambda) \,.
\]
The sine-process null's second-moment budget is the unconditional prime-side
$\kappa(\lambda)$ of Alp\"oge and Furman, exactly, on the whole band
$\lambda \le 1$ --- an exact identity check at second order; the $F'$ budget row
is therefore backed by unconditional data. (Provenance: the prime-side member
--- $\tr \Gh^2 \leftrightarrow \int (\lambda - |\alpha|) F(\alpha)\, d\alpha$
with $F$ unconditional and pointwise on the closed band $0 \le \alpha \le 1$
--- is the Theorem~5.8 / Remark~5.10 bookkeeping of~\cite{AF26} together with
Theorem~1 of~\cite{BGSTB24}, quoted, not re-derived; the displayed arithmetic
identity is the elementary integral
$\int_0^\lambda (\lambda - a) a\, da = \lambda^3/6$, doubled, plus the delta
term $\lambda$.) No third-order analog is claimed: the identification
$c_3(\lambda') = m_3(\lambda')$ rests on Rudnick--Sarnak-range GUE $3$-level
correlations and remains open (the cubic budget center;
Section~\ref{sec:flags}).
\end{corollary}

Numerical confirmation (\texttt{verify\_t4.py}, \texttt{verify\_t4\_mc.py}):
Fourier-side quadrature matches both closed forms at $7$ bandwidths to
$\le 1.7 \times 10^{-5}$; an independent real-space determinantal quadrature
gives $m_3(1/2) = 4.9924$ (truncation $X = 40$) $\to 4.9960$ ($X = 80$),
trending to $5$ as a second independent quadrature did ($4.9896 \to 4.9932$); a
third independent sampler, CUE Monte Carlo, extrapolates $m_2(1/2) \to 2.1668$
and $m_3(1/2) \to 5.0014$ against $13/6$ and $5$, each within ${\sim}1\sigma$ of
the other two (full tables: Appendix~\ref{app:mc}). The values
$m_2(1) = 4/3$, $m_3(1) = 2$ also discharge, for $k \le 3$, the sine moments
$m_k(1) = 1, 4/3, 2$ recalled in~\cite{AF26}.

\subsection{Grid Parseval decoupling}
\label{sec:decoupling}

The absorption mechanism rests on one exact identity: on a specific uniform grid
the entire bandwidth-one Frobenius row degenerates to the multiplicity count ---
for every site subset and every mark assignment --- while the half-band kernel on
the same grid remains position-sensitive.

\begin{lemma}[flat-band lemma]
\label{lem:flatband}
Let $M \ge 1$ and let the harmonic band be any set of $M$ consecutive integers
$B = \{j_0, \ldots, j_0 + M - 1\}$, with uniform weights $u_j = 1/M$. Let the
configuration be supported on the $M$-site uniform grid $\theta_k = k N/M$
($k = 0, \ldots, M - 1$) on the circle of circumference $N$, with arbitrary real
marks $m_k \ge 0$. Then
\[
  \tr \Gh^2 = \sum_k m_k^2 \qquad \text{exactly.}
\]
\end{lemma}

\begin{proof}
On the grid the form factor is a DFT on $\Z_M$:
$c_s = \sum_k m_k e^{-2\pi i s k/M}$, which depends only on $s \bmod M$ --- $c$
is $M$-periodic in $s$ --- and DFT Parseval gives
$\sum_{r \bmod M} |c_r|^2 = M \sum_k m_k^2$ (P). The Frobenius row is
$\sum_s W_2(s) |c_s|^2$ with
$W_2(s) = \#\{(j_1, j_2) \in B \times B : j_1 + j_2 = s\}/M^2$; the sum-set
$B + B = \{2 j_0, \ldots, 2 j_0 + 2M - 2\}$ carries the triangular pair counts
$M - |s - (2 j_0 + M - 1)|$ on $|s - (2 j_0 + M - 1)| \le M - 1$. Fix a residue
class $r \bmod M$: the window $B + B$ has length $2M - 1$, so it contains either
one member of the class ($t = 0$, count $M$) or two members $t$ and $t \mp M$
with $0 < |t| < M$, whose counts sum to $(M - |t|) + |t| = M$ --- exactly $M$ on
every residue class. Hence, by the $M$-periodicity of $|c_s|^2$,
\[
  \tr \Gh^2 = \frac{1}{M^2} \sum_{r \bmod M} M |c_r|^2
  = \frac{1}{M} \sum_r |c_r|^2 = \sum_k m_k^2
\]
by (P).
\end{proof}

\begin{theorem}[grid Parseval decoupling at bandwidth one]
\label{thm:parseval}
Let $N$ be even, flat window at $\lambda = 1$: band $B = \{|j| \le N/2\}$,
$M = N + 1$ harmonics, $u_j = 1/(N+1)$. Let the configuration be supported on the
$(N+1)$-site uniform grid of spacing $N/(N+1)$ (the $\psi_1$-zero grid), with
arbitrary marks. Then
\[
  F_1 = \tr \Gh_1^2 = \sum_z m_z^2 \qquad \text{exactly,}
\]
for every site subset and every mark assignment: the bandwidth-one reading on
this grid sees only multiplicities, never positions, and the
\texttt{lemmaR\_tight} extremal analysis transfers verbatim to grid
configurations.
\end{theorem}

\begin{proof}
Lemma~\ref{lem:flatband} with $M = N + 1$, $B = \{-N/2, \ldots, N/2\}$ ($N + 1$
consecutive integers since $N$ is even), grid spacing $N/M = N/(N+1)$.
\end{proof}

\paragraph{Nomenclature.}
Lemma~\ref{lem:flatband} is critical-sampling Parseval for a flat band --- the
discrete orthogonality of $M$ consecutive characters on an $M$-site uniform grid
--- and Proposition~\ref{prop:lattice} is its Shannon--Whittaker continuum form
on the integer lattice. Both are classical; they are proved here in the exact
finite-circle normalization of Section~\ref{sec:finitecircle}, and the Lean
development machine-checks Lemma~\ref{lem:flatband} over an arbitrary integral
domain (Section~10). The degeneracy they create at critical sampling is the one
Lagarias and Rodgers study in~\cite{LR21}: above the Nyquist bandwidth the
mimicking correlations are unique, at critical sampling they are not. The grid
used here (spacing $a = N/(N+1)$, band $B = 1/2$) falls in the region their
Theorem~1.8 leaves undetermined, and the witness of
Section~\ref{sec:certification} matches finitely many trace rows rather than all
correlations, so there is no conflict; their negative results are why the
absorption cannot be upgraded to full $n$-level mimicry.

\begin{remk}[why this grid]
Equivalently: all pairwise site differences are nonzero multiples of $N/(N+1)$,
the zero set of the discrete bandwidth-one kernel $D$. The Fourier proof makes
the exactness structural --- assembly weights constant on residue classes ---
rather than a numerical coincidence.
\end{remk}

\begin{proposition}[the $\lambda' = 1/2$ kernel is alive on the same grid]
\label{prop:halfband}
Same geometry, $N$ and $N/2$ even, $\lambda' = 1/2$: band
$B' = \{|j| \le N/4\}$, $M' = N/2 + 1$. For grid-supported configurations,
\[
  F' = \sum_{|s| \le N/2} \frac{M' - |s|}{M'^2}\, |c_s|^2 \,,
\]
where $\{-N/2, \ldots, N/2\}$ is a complete residue system mod $(N+1)$ and $c_s$
is the $\Z_{N+1}$ DFT of the mark vector. The weights are strictly positive and
non-constant across residues (triangular: at $N = 64$, from $1/33$ at $s = 0$
down to $1/33^2$ at $|s| = 32$), so $F'$ is a strictly position-sensitive
functional of the grid configuration, subject only to Parseval.
\end{proposition}

\begin{proof}
$F' = \sum_s W_2'(s) |c_s|^2$ with $W_2'(s) = (M' - |s|)_+/M'^2$ supported in
$|s| \le N/2$. That range has $N + 1$ elements, hence meets each residue class
mod $N + 1$ exactly once; no residue-class telescoping occurs (band length
$M' = N/2 + 1$ differs from the grid size $N + 1$), so $F'$ reads the full grid
power spectrum with a non-uniform weight; distinct site subsets of equal mark
statistics generically differ.
\end{proof}

\paragraph{Worked exact value: the witness's first column.}
$N = 64$, the ``vacancy lattice'': $64$ simple marks on $64$ of the $65$ grid
sites. Then $c_0 = 64$ and $|c_s|^2 = 1$ for $s \not\equiv 0 \bmod 65$, so
\begin{align*}
  F_1 &= 64 \qquad (\text{Theorem~\ref{thm:parseval}}),\\
  F' &= \frac{33}{33^2}\, 64^2
    + \sum_{0 < |s| \le 32} \frac{33 - |s|}{33^2}
    = \frac{4096}{33} + \frac{32}{33}
    = \frac{4128}{33} = 125.0909\ldots \,,
\end{align*}
using $\sum_{|s| \le 32} (33 - |s|) = 33^2$ and subtracting the $s = 0$ term.
This reproduces the certified witness column exactly
(\texttt{witness\_N64.json}: $F_1 = 64$, $F' = 125.09090\ldots$). The
two-bandwidth decoupling is thereby explicit: $F_1$ is blind to the vacancy's
position, $F'$ is not.

\begin{proposition}[continuum analog: the adversary lattice]
\label{prop:lattice}
In the continuum model, flat window, configuration supported on the integer
lattice $\Z$ with arbitrary real marks: $\tr G_1^2 = \sum m_z^2$ exactly, while
the half-band kernel is alive on the lattice:
$\psi_{1/2}(k) = \sinc(\pi k/2) = (2/(\pi k))(-1)^{(k-1)/2}$ for odd $k$ ($0$ for
even $k \ne 0$). For the finite-$T$ Weil--Gabor matrix itself the identity holds
as $\tr \Gh^2 = (\sum m^2)(1 + o(1))$, the $o(1)$ being the finite-window
corrections of~\cite{AF26} (quoted).
\end{proposition}

\begin{proof}[Proof (model part)]
At $k = 2$ the master formula gives
$\tr G_1^2 = \sum_{z, z'} m_z m_{z'} \psi_1(z - z')^2$;
$\psi_1(x) = \sinc(\pi x)$ vanishes at every nonzero integer and equals $1$ at
$0$, so only the diagonal survives. The half-band values are
$\sin(\pi k/2)/(\pi k/2)$ with $\sin(\pi k/2) = (-1)^{(k-1)/2}$ for odd $k$.
\end{proof}

\begin{remk}
Theorem~\ref{thm:parseval} is the exact periodization of
Proposition~\ref{prop:lattice}: spacing $N/(N+1)$ instead of $1$ makes the
identity exact rather than asymptotic, which is why the finite-$N$ computation
exhibits $\delta_0 = 0$ as an exact LP equality rather than a numerical
near-zero. The absorption is not a finite-$N$ artifact.
\end{remk}

Numerical confirmation (\texttt{verify\_t1.py}): grid exactness on $40$ random
site-subsets with random marks at $N = 64$ and $N = 128$, two independent
evaluation routes (Fourier double sum; eigenvalues of $H$):
$\max |F_1 - \sum m^2| = 9.1 \times 10^{-13} / 4.1 \times 10^{-12}$, matching
the independently computed $4.5 \times 10^{-13}$; general flat-band lemma on
$20$ random $(N, M, j_0)$ draws including asymmetric bands:
$3.2 \times 10^{-12}$; controls --- off-grid configurations give
$|F_1 - \sum m^2| \in [1.3 \times 10^{-2},\ 62.5]$ and on-grid half-band
$|F' - \sum m^2| \in [40.9, 120.8]$ (the identity genuinely fails off the grid,
and the half-band kernel is genuinely alive on it); continuum lattice exact to
machine precision.

\section{The absorption theorem}
\label{sec:absorption}

This section states and proves the model-level theorem the LP pins --- the
two-moment data force exactly the \texttt{lemmaR\_tight} corner, for both
benchmarks, over the model's configuration classes --- and then gives the
computational certification that the entire cubic block is absorbed at that
corner, with the explicit witness.

\subsection{The atom-only corner theorem}
\label{sec:atomonly}

Fix $N$ and $\varepsilon \ge 0$, and let $\mathcal{A}$ be the class of atom-only
configurations: atoms at arbitrary positions (finite-circle or continuum model,
flat windows --- the proofs work verbatim in both), integer marks $m_i \ge 1$,
mass $\sum m_i = N$ per period.

\begin{lemma}[diagonal positivity]
\label{lem:diagpos}
For every atom-only configuration (any positions, any real marks $m_i \ge 0$), in
either model, $F_1 \ge \sum_i m_i^2$, with equality iff every off-diagonal kernel
value vanishes --- in particular on the grid of Theorem~\ref{thm:parseval} and
the lattice of Proposition~\ref{prop:lattice}.
\end{lemma}

\begin{proof}
$F_1 = \sum_{i, i'} m_i m_{i'} K(\theta_i - \theta_{i'})$ with $K = D^2$ (finite
circle) or $K = \psi_1^2$ (continuum): $K \ge 0$ pointwise, $K(0) = 1$, so
$F_1 = \sum m_i^2 + (\text{a sum of nonnegative off-diagonal terms})$.
\end{proof}

This is the one step that fails for pair-containing configurations --- an
off-line pair contributes cross-terms of the form
$2 m m' \operatorname{Re}[\psi(x + iy)^2]$, which need not be nonnegative;
Sections~\ref{sec:fractional}--\ref{sec:pairchannel} treat that channel.

\begin{lemma}[integer-mark corner inequalities]
\label{lem:cornerineq}
For every atom-only configuration with integer marks $m_i \ge 1$ and mass $N$:
\begin{align*}
  \text{(a)} \quad &\Nd = \#\mathrm{atoms}
    \ \ge\ \frac{3N - \sum_i m_i^2}{2}\,,\\
  \text{(b)} \quad &n_1 = \#\text{mark-}1\ \mathrm{atoms}
    \ \ge\ 2N - \sum_i m_i^2 \,,
\end{align*}
each with equality iff all marks lie in $\{1, 2\}$.
\end{lemma}

\begin{proof}
Per-atom inequalities summed over atoms. (a) For integer $m \ge 1$,
$(m - 1)(m - 2) \ge 0$, i.e.\ $1 \ge (3m - m^2)/2$, equality iff
$m \in \{1, 2\}$; sum. (b) $\one_{m = 1} \ge 2m - m^2$: equality at $m = 1, 2$;
at $m \ge 3$ the right side is negative; sum.
\end{proof}

\begin{remk}
(a) is the second integrality level $(m-1)(m-2) \ge 0$ --- exactly the level
\texttt{lemmaR\_tight} proves exhausted by the two-moment certificate, and, by
Theorem~\ref{thm:blind}, exactly the level the cubic row re-measures on isolated
blocks; (b) is its simple-fraction sibling. The device is Conrey, Ghosh and
Gonek's~\cite{CGG98}, made unconditional in \S7.2(c) of~\cite{AF26}. In the
terms of that paper, (a) is the doubles-optimal corner arithmetic: an atom
of mark $m$ spends $m(m-1)$ of Frobenius excess to save $m - 1$ units of
$\Nd$, so doubles are the efficient coin.
\end{remk}

\begin{theorem}[atom-only corner theorem, both benchmarks]
\label{thm:corner}
Fix $N$ even, $0 \le \varepsilon \le 1/2$ (finite-circle model, flat windows).
Over laws $w$ on $\mathcal{A}$ with mass $N$ and
$E_w[F_1] \le (4/3) N (1 + \varepsilon)$:
\begin{align*}
  \text{(a)} \quad &\min_w\ E_w[\Nd]/N = \frac{5}{6} - \frac{2}{3}\varepsilon \,,\\
  \text{(b)} \quad &\min_w\ E_w[p_1] = 2 - \frac{4}{3}(1 + \varepsilon)
    = \frac{2}{3} - \frac{4}{3}\varepsilon \,.
\end{align*}
Both minima are attained by explicit laws supported on the $(N+1)$-site grid of
Theorem~\ref{thm:parseval} with marks $\{1, 2\}$ only. Every configuration in
$\mathcal{A}$ satisfies the pointwise versions with its own Frobenius excess
$\varepsilon_c$. The same statements hold with the two-sided budget band
$(4/3) N [1 - \varepsilon, 1 + \varepsilon]$ (the attaining laws sit at the upper
edge), and the minima over any richer mark alphabet (any $W$, or unbounded marks)
are the same.
\end{theorem}

\begin{proof}
Lower bounds: by Lemma~\ref{lem:diagpos},
$E_w[\sum m^2] \le E_w[F_1] \le (4/3) N (1 + \varepsilon)$; by
Lemma~\ref{lem:cornerineq}(a) and linearity,
$E_w[\Nd] \ge (3N - (4/3)N(1 + \varepsilon))/2
 = N(5/6 - (2/3)\varepsilon)$; by Lemma~\ref{lem:cornerineq}(b),
$E_w[n_1] \ge 2N - (4/3)N(1 + \varepsilon) = N(2/3 - (4/3)\varepsilon)$.

Attainment: for integer $0 \le k \le N/2$ let $c_k$ be the grid configuration
with $k$ doubles and $N - 2k$ simples on $N - k$ distinct grid sites (which
exist: $N + 1 \ge N - k$). By Theorem~\ref{thm:parseval},
$F_1(c_k) = N + 2k$ exactly; $\Nd(c_k) = N - k$; $n_1(c_k) = N - 2k$. Set
$\bar{k} = (B - N)/2$ with $B = (4/3)N(1 + \varepsilon)$; then
$0 \le \bar{k} \le N/2$ for $\varepsilon \le 1/2$. The two-column law $w^*$
mixing $c_{\lfloor \bar{k} \rfloor}$ and $c_{\lfloor \bar{k} \rfloor + 1}$ with
$E[k] = \bar{k}$ achieves
\begin{gather*}
  E[F_1] = N + 2\bar{k} = B \quad (\text{upper edge, exactly}),\\
  \frac{E[\Nd]}{N} = \frac{3N - B}{2N} = \frac{5}{6} - \frac{2}{3}\varepsilon,
  \qquad
  E[p_1] = 2 - \frac{B}{N} = \frac{2}{3} - \frac{4}{3}\varepsilon \,,
\end{gather*}
and both lower-bound chains hold with equality (marks in $\{1, 2\}$; grid kills
all cross-terms; budget saturated), so $w^*$ is simultaneously optimal for both
objectives. The two-sided band adds only the lower budget edge, which $w^*$
satisfies. Enlarging the mark alphabet cannot lower the minimum (the chain never
used a mark cap) and cannot raise it ($w^*$ uses marks $\{1, 2\}$, inside every
alphabet): the $W$-independence is inclusion, not extrapolation.
\end{proof}

\begin{corollary}[asymptotic and matched corners]
\label{cor:matched}
As $\varepsilon \to 0$ the two optima tend to $5/6$ and $2/3$ --- the model's
bandwidth-one analogs of Montgomery's extremal values~\cite{Mon73} for distinct
zeros and simple fraction. With the matched-geometry budget $B_1$ replacing
$(4/3)N$, the same proof gives
$\min E[\Nd]/N = (3N - B_1(1 + \varepsilon))/(2N)$ and
$\min E[p_1] = 2 - (B_1/N)(1 + \varepsilon)$; at the matched values used here
($N = 64$, $\varepsilon = 0.05$, $B_1 = 84.71214548308222$) these are
$0.8050957$ and $0.6101914$ --- the computed optima to LP
precision.
\end{corollary}

\begin{remk}[a numerical coincidence, disambiguated]
The corner value $0.8050957$ rounds to $0.8051$ at four decimals, and the string
$0.8051$ also occurs in the zeta literature as Farmer--Gonek--Lee's RH-conditional
pair-correlation lower bound for the proportion of \emph{distinct} zeros of
zeta,
$\Nd(T) > 0.8051\,N(T)$ \cite[Cor.~1.4]{FGL14} (since superseded by $0.84665$
in~\cite{BHB13} and by $0.8477$ on RH / $0.8486$ on GRH
in~\cite[Cor.~3]{CGdL20}). The two are unrelated: the value here is a
finite-$N$ LP
corner of the bandwidth-one model at $N = 64$, $\varepsilon = 0.05$,
$B_1 = 84.71214548308222$, and carries no information about zeta.
\end{remk}

Numerical confirmation (\texttt{verify\_t2.py}): per-atom inequalities exhaustive
to $m = 50$ plus $20{,}000$ random draws ($0$ violations); grid-column LP matches
the closed forms at $\varepsilon \in \{0.10, 0.05, 0.02, 0.002\}$ to
$\le 3.3 \times 10^{-16}$ for both objectives; exact rational identities verified
in \texttt{Fraction} arithmetic; the certified witness reproduces
$P = (3N - E[F_1])/(2N) = 0.8050957$ to $6 \times 10^{-13}$.

\subsection{Integrality is necessary: the exact fractional counterexample}
\label{sec:fractional}

Any extension of Theorem~\ref{thm:corner} beyond the atom-only class must
consume integrality, as the following counterexample shows. Work in the primary
geometry ($N = 64$, flat window, $\lambda = 1$; kernel
$\psi(z) = \sum_j u_j e^{-2\pi i j z/N}$, grid $g_k = k\,64/65$) and define
$\bar{a}(x) := \psi(ix) = \sum_j u_j \cosh(2\pi j x/N) \ge 1$. Write $M(c)$ for
total mark mass, $S_2(c) = \sum_z m_z^2$ (each pair member counted $m^2$),
$T(c) = 3M(c) - 2\Nd(c)$, and consider the \emph{master inequality}
\[
  \text{(MI)} \qquad F_1(c) \ \ge\ 3M(c) - 2\Nd(c) \,,
\]
which is precisely what the corner consumes: (MI) per configuration gives, for
every law with mass $N$ and $E[F_1] \le (4/3)N(1 + \varepsilon)$, both corners ---
$E[\Nd]/N \ge 5/6 - (2/3)\varepsilon$ and
$E[p_1] \ge 2/3 - (4/3)\varepsilon$. (Per configuration (MI) reads
$\Nd \ge (3M - F_1)/2$; take expectations. For simples,
$\one_{m = 1} \ge 2 - m$ summed over the zero multiset gives
$n_1 \ge 2\Nd - M$, hence $E[n_1] \ge (3N - E[F_1]) - N$.) Note $S_2 \ge T$ for
integer marks (per zero, $m^2 - (3m - 2) = (m-1)(m-2) \ge 0$), so (MI) is weaker
than $F_1 \ge S_2$; the closure arguments below target (MI) because it carries
the integrality slack explicitly.

\begin{proposition}[fractional marks break the floor at every depth]
\label{prop:fractional}
Let $c_\mu$ be the vacancy lattice (unit atoms on grid sites
$g_1, \ldots, g_{64}$) plus one pair at the hole $g_0$ with depth $d > 0$ and
real mark $\mu > 0$. Then, exactly,
\[
  F_1(c_\mu) - S_2(c_\mu)
  = 2\mu^2 \bar{a}(2d)^2 - 4\mu\,(\bar{a}(d)^2 - 1),
\]
which is negative for every $0 < \mu < 2(\bar{a}(d)^2 - 1)/\bar{a}(2d)^2$ --- the
pair harvests the hole's missing coupling linearly in $\mu$ while paying only
quadratically. For integer marks the same family is safe: at $\mu = m \in \Z$,
the Cauchy--Schwarz bound $\bar{a}(d)^2 - 1 \le (\bar{a}(2d) - 1)/2$ gives
$F_1 - S_2 \ge 2m^2 \bar{a}(2d)^2 - 2m(\bar{a}(2d) - 1) > 0$.
\end{proposition}

\begin{proof}
On the vacancy lattice the atom form factor is $\varphi_0(s) = 64$ at
$s \equiv 0$ and $-1$ otherwise; the pair adds $2\mu \cosh(\beta s)$,
$\beta = 2\pi d/N$. Expanding $F_1 = \sum_s w_s |c_s|^2$ with the generating
identities $\sum_s w_s \cosh(\beta s) = \bar{a}(d)^2$ and
$\sum_s w_s \cosh(2\beta s) = \bar{a}(2d)^2$ (the assembly weights $w = u * u$
factor through $\psi$ at imaginary arguments) gives the display; the numeric
check at $(d, \mu) = (0.25, 0.05)$ reproduces it to $10^{-12}$, with
$F_1 - S_2 = -3.520 \times 10^{-2}$. The integer case uses Cauchy--Schwarz on the
positive weights $u_j$, which sum to $1$:
$\bar{a}(d)^2 = \bigl(\sum_j u_j \cosh(2\pi j d/N)\bigr)^2
 \le \sum_j u_j \cosh^2(2\pi j d/N) = (1 + \bar{a}(2d))/2$, the last equality
by $\cosh^2 t = (1 + \cosh 2t)/2$.
\end{proof}

\paragraph{Reading.}
An adversary with fractional marks takes $\mu \to 0$ and wins at any depth;
integer marks force $\mu \ge 1$, where the quadratic cost dominates. Any closure
of the pair channel must therefore consume integrality --- the channel is
invisible to every convexified or soft-positivity analysis. (This is the textbook
\emph{integrality gap} in its usual form: the fractional relaxation of the
marked-configuration program has a strictly smaller optimum than the integer
program, so a bound proved for the relaxation is not a bound for the problem.)
It also explains the LP's phenomenology: its columns are explicit integer
configurations (nothing about marks is relaxed), which is why its converged
pricing found no pair column. It is likewise the sharp form of the realizability
discussion of Section~\ref{sec:adversary}: with fractional marks even the $5/6$
baseline is false.

\subsection[The pair-interference channel: closed on the pair-depth family
(multi-pair to w <= 0.98), backstopped everywhere]{The pair-interference channel:
closed on the pair-depth family (multi-pair to $w \le 0.98$), backstopped
everywhere}
\label{sec:pairchannel}

The full closure argument, with its capacity machinery and certified constants,
is given in the companion note
\texttt{results/a4-no-go/}\allowbreak\texttt{pair-}\allowbreak\texttt{channel.md}
of the repository (see \emph{Provenance, data and code} at the end of this
paper). This subsection states the results, the exact ledger they rest on, and
the proofs that are short enough to give in full. Everything is at the
primary geometry ($N = 64$, flat window, $\lambda = 1$); depths $d$ are in
circle units, $w = 2\pi d$ dimensionless (the admissible pair family of
Section~\ref{sec:adversary} is
$w \in (0, 1]$ plus deep probes $w \in \{1.5, 2\}$, i.e.\ $d \le 0.3183$).

\paragraph{The exact interference ledger.}
For every admissible configuration (any atoms, any pairs, any integer marks),
with $R_y(x) := \operatorname{Re} \psi(x + iy)^2$:
\begin{align*}
  F_1 - T
  &= \sum_a (m_a - 1)(m_a - 2) + 2 \sum_p (m_p - 1)(m_p - 2)
    && [\text{integrality slack}]\\
  &\quad + \sum_{a \ne a'} m_a m_{a'} \psi(x_{aa'})^2
    && [\text{atom-atom} \ge 0]\\
  &\quad + \sum_p 2 m_p^2 \bar{a}(2 d_p)^2
    && [\text{pair diagonals}]\\
  &\quad + \sum_p 4 m_p \sum_a m_a R_{d_p}(x_{pa})
    && [\text{pair-atom}]\\
  &\quad + \sum_{p < q} 4 m_p m_q
      \bigl[ R_{d_p + d_q}(x_{pq}) + R_{|d_p - d_q|}(x_{pq}) \bigr]
    && [\text{pair-pair}]
\end{align*}
(verified on $25$ random mixed configurations to $2.2 \times 10^{-11}$, the
stored \texttt{pairchan\_verify\_out.json} maximum). Everything on the right is
nonnegative except the $R$-couplings, whose negative parts live only in the dips
of $R_y$; the question is whether the dips can outrun the pair diagonals plus the
integrality slack. They cannot: at the worst depth the margin is a factor of
about $1.5$.

\begin{theorem}[single-pair closure]
\label{thm:singlepair}
Every configuration with exactly one pair (arbitrary atoms, arbitrary integer
marks, arbitrary positions) satisfies (MI) whenever the pair depth is
$d \le 0.45$ --- i.e.\ $w \le 2.8$, covering the entire admissible pair-depth
family
including both deep probes, with $\ge 34\%$ capacity margin on the family itself
(certified stacking-corrected capacity ratio $\le 0.66$ on the family,
$\le 0.921$ out to $d = 0.45$). For mass $\le N$ configurations (the LP class)
(MI) also holds for every $d \ge 1.0$; the window $(0.45, 1.0)$ is certified
numerically (raw capacity ratio $\le 0.6823$ throughout; direct adversarial
searches reach at most $74\%$ of the pair budget).
\end{theorem}

\begin{theorem}[multi-pair closure]
\label{thm:multipair}
(MI) holds for every configuration all of whose pair depths satisfy
$w \le 0.82$, unconditionally; and for all depths $w \le 0.98$ modulo one
isolated cell-crowding cap whose ledger constant is interval-certified
($\sup \Sgen \le 0.98465 < 1$ on $(0, 0.156]^2$, per-cell rigorous
enclosures). On the deepest sliver $w \in (0.98, 1]$ the ledger inequality fails
(interval-certified $\Sgen \ge 1.01405 > 1$ at $d = d' = 0.159$); there the
targeted crowd-plus-sea attacks --- $8$ to $65$ pairs, one per cell tranche, plus
greedy mark-$2$ seas up to $130$ atoms, re-run at $d = 0.158$ and $0.159$ ---
never bring $F_1 - T$ below $+17.46$, against a floor requirement of $0$, and the
unconditional $8/9$ backstop applies.
\end{theorem}

\begin{proof}[Proof architecture (full details in the companion note)]
\renewcommand{\qedsymbol}{QED (architecture)}
Partition the circle into $65$ kernel-zero cells; per cell, the exploitable dip
mass $\kappa_k(y) = \sup_{\mathrm{cell}} [-R_y]_+$ sums to the capacity curve
$\nu(y)$, with the exact curvature identity
$\alpha := \sum_j u_j (2\pi j/N)^2 = \sum_k \psi'(g_k)^2 = 3.392677$ controlling
the shallow regime ($\nu(y) = \alpha y^2 (1 + o(1))$). Integer atoms harvesting a
cell's dip pay in-cell atom-atom crosses ($\ge \psi(\text{dip-zone diameter})^2$
each) and their own integrality slack, capping the harvest at
$8 m_p \kappa_k$ per cell for marks $\le 2$ (with a certified stacking correction
beyond); for pair depths $w \le 1$ the closing inequality
$8\nu(y) \le 2\bar{a}(2y)^2$ follows in closed form from one harmonic-Taylor
bound with the exact $\alpha$ and one certified finite constant ($1.2\%$ closing
margin at $w = 1$); the deep probes are covered by the certified capacity ratios
of Theorem~\ref{thm:singlepair}. Pair-pair couplings are priced through the
joint kernel
$\Phi(d, d'; x) = R_{d+d'}(x) + R_{|d-d'|}(x)$ (never the two parts separately),
writing $\nu_{\mathrm{joint}}(d, d') = \sum_k \sup_{\mathrm{cell}} [-\Phi]_+$ for
its capacity curve and
$\Phi_0 = \inf\{ R_s(x) + R_t(x) : |x| \le 32/65,\ s \le 1/\pi,\ t \le 1/(2\pi) \}$
for the in-cell coupling a second pair must pay; the ledger constant is then
$\Sgen(d, d')
 = [8\nu(d) + 4\nu_{\mathrm{joint}}(d, d')] / (2\bar{a}(2d)^2)$ and the closure
needs $\Sgen \le 1$. This carries a
one-pair-per-cell crowding cap that is self-consistent for $w \le 0.82$ (both
sides interval-certified:
$2 \max \nu_{\mathrm{joint}} \le 0.60910 < 0.64371 \le \Phi_0$) and whose
ledger constant beyond it is interval-certified to
$w \le 0.98$; beyond $w = 0.98$ the ledger fails (certified) and the closure
defers to numerics plus the backstop. Deep single pairs ($d \ge 1$) fall to the
crude bound $\bar{a}(2d)/\bar{a}(d) \ge 12.158$ nondecreasing.
\end{proof}

\begin{theorem}[equal-depth pair-only configurations, exact, all depths]
\label{thm:equaldepth}
If $c$ consists only of pairs, all at one common depth $d$ (any positions, any
integer multiplicities), then $F_1 \ge S_2 + 2 \sum_p m_p^2$ --- (MI) with a full
doubling-cushion surplus at every depth.
\end{theorem}

\begin{proof}
$c_s = 2\cosh(2\pi s d/N)\,\varphi(s)$ with $\varphi$ the position-multiset form
factor, and $4\cosh^2 = 2 + 2\cosh(\text{double})$: $F_1$ splits as twice the
atom-only Frobenius of the position multiset ($\ge 2\sum_p m_p^2$ by diagonal
positivity) plus a termwise-$\cosh$-weighted copy ($\ge$ the same floor).
\end{proof}

\begin{theorem}[unconditional spectral backstop, no restrictions]
\label{thm:backstop}
For every admissible configuration --- any number of pairs, any depths, any
marks --- $\Nd \ge (4/9)(3M - F_1)$. Hence for any law with mass $N$ and
$E[F_1] \le (4/3)N(1 + \varepsilon)$:
\[
  \frac{E[\Nd]}{N} \ \ge\ \frac{8}{9}\Bigl( \frac{5}{6}
    - \frac{2}{3}\varepsilon \Bigr)
  = \frac{20}{27} - \frac{16}{27}\varepsilon
  = 0.7407\ldots - 0.593\,\varepsilon \,,
\]
so the maximal conceivable pair advantage on the corner, in any regime whatever,
is $\le 5/6 - 20/27 = 5/54 = 0.0926$, unconditionally.
\end{theorem}

\begin{proof}
Let $B$ be the Hermitian frequency matrix
$B[j, j'] = \sqrt{u_j u_{j'}}\, c_{j - j'}$ ($65 \times 65$): exactly
$\tr B = M$ and $\tr B^2 = F_1$. $B$ decomposes as
$\sum_{\mathrm{atoms}} m_a v_a v_a^{\dagger} + \sum_{\mathrm{pairs}} B_p$ with
each atom block PSD of rank $1$. Each pair block is
$B_p = m(p q^{\dagger} + q p^{\dagger})$, with
$p_j = \sqrt{u_j}\, e^{-2\pi i j \theta/N} e^{2\pi j d/N}$ and
$q_j = \sqrt{u_j}\, e^{-2\pi i j \theta/N} e^{-2\pi j d/N}$ (the two summands of
the conjugate-pair form factor of Section~\ref{sec:finitecircle}); on
$\operatorname{span}\{p, q\}$ its nonzero eigenvalues are
$m(\operatorname{Re}\langle q, p\rangle \pm |p||q|)$, and Cauchy--Schwarz gives
$\operatorname{Re}\langle q, p\rangle - |p||q| \le 0$, so $n_+(B_p) \le 1$ ---
the finite-circle counterpart of the signature-$(1,1)$ block of
Proposition~\ref{prop:pairblock}(iv). By subadditivity of positive inertia,
$n_+(B) \le \#\mathrm{atoms} + \#\mathrm{pairs} \le \Nd$. For every real $t$,
$3t - t^2 \le (9/4)\one_{t > 0}$; summing over the spectrum,
$3M - F_1 = \sum_i (3\lambda_i - \lambda_i^2) \le (9/4) n_+(B) \le (9/4)\Nd$.
\end{proof}

\begin{remk}
The $9/4$ (rather than $2$) is the price of a purely spectral argument ---
$3t - t^2$ exceeds $2$ on $(1, 2)$ with max $9/4$ at $t = 3/2$ --- the same
structural reason the inertia route of \cite{AF26} cannot reach $5/6$ by itself.
The bound is, in compensation, depth-uniform, mark-uniform and
interference-proof.
\end{remk}

\paragraph{Status and the sharp conjecture.}
Roughly $10^4$ adversarially optimized configurations (dip attacks, pair crowds,
half-gap frustrated twins, deep seas, random mixed stress; seeds fixed,
independent assembly) never brought $F_1 - T$ below $+2.000000$, and the minimum
of $F_1 - T - \sum_p 2(2 m_p^2 - 3 m_p + 2)$ over every family is $0.000000$,
attained only in the depth $\to 0$ grid limit. The searches establish that the
binding adversary is always shallow; that coordinated dip attacks can consume
nearly the whole depth-dependent excess $2m^2(\bar{a}(2d)^2 - 1)$ (within $3\%$
at $d = 0.1$) but never any of the depth-independent doubling cushion $2m^2$ ---
exactly the split the capacity lemmas predict. We record the sharp conjecture
(numerically exhaustive, not proved):
$F_1 \ge \sum_a (3 m_a - 2) + 4\sum_p m_p^2$ --- every pair pays its full
flattened-double cost and interference recovers nothing. Its near-tight witnesses
($+0.03$) show any proof must be capacity-sharp, which is why the theorems above
target (MI), with its floor headroom of $2$ per pair, instead.

\paragraph{Consequence.}
The model-level corner theorem extends from the atom-only class to the bulk of
the admissible class of Section~\ref{sec:adversary}: single-pair columns at
every admissible depth ($w \le 2.8$), multi-pair columns at $w \le 0.82$
unconditionally and
$w \le 0.98$ modulo the crowding cap (constant interval-certified), equal-depth
pair-only columns at every depth --- for both benchmarks, via (MI). The remaining
admissible regimes --- multi-pair columns at the deep-probe depths
$w \in \{1.5, 2\}$ (inside the admissible class but beyond the theorems' range),
multi-pair columns with a depth in $w \in (0.98, 1]$ (ledger failure certified
there; coverage is the re-run attacks plus the backstop), and mass-unbounded
crowds (not LP-relevant, since columns have mass $N$ and deeper objects are
garnish priced by Section~\ref{sec:capacity}'s ladder) --- are floored at $8/9$
of the corner unconditionally and covered by converged pricing (O2); the
$(0.98, 1]$ sliver's coverage is the re-run attacks under (O1) plus the backstop.
Class-level exactness of $\delta_0 = 0$ is therefore proved on the closure's
regimes and heuristic (pricing-backed, backstopped at $20/27$) on that residue.
A pair-assisted violation would have registered as a strictly positive marginal
value, to be cut by pricing rather than by absorption; none is present on the
closure's regimes. The residual coverage is stated in Section~\ref{sec:flags}.

\subsection[The computational certification: 940 records, the epsilon law, and
the witness]{The computational certification: 940 records, the $\varepsilon$
law, and the witness}
\label{sec:certification}

The decision LP (column generation over explicit configurations, analytic-gradient
multi-start pricing, exact-LP equality on the accumulated dictionary, final
solutions re-verified in \texttt{Fraction} arithmetic on $10^{-12}$
rationalizations) was run over the full decision grid. The outcome at every
record is absorption; the records are stored in machine-readable form.

\paragraph{The grid.}
$N = 64$: $\{$matched, asymptotic$\}$ budgets $\times$
$W \in \{2, 3, 4, 6, 8\}$ $\times$ $\Cled \in \{10, 100, 1000\}$ $\times$
$\varepsilon \in \{0.02, 0.05, 0.10\}$ $\times$ fuzz $\{$coupled, none,
$\Gamma \in \{0.05, 0.10, 0.25\}\}$ $\times$ budget setting $\{$centered, set
adversary-favorably within the error bars$\}$ $= 900$
records. $N = 128$: $\{$matched, asymptotic$\}$ $\times$ $W$ $\times$ fuzz
$\{$coupled, none$\}$ $\times$ the same two budget settings $= 40$ records.
Total 940 records; every record solved to status success; LP duality gap
${\sim}10^{-9}$.

\paragraph{The result.}
$\delta_0 = P_{\mathrm{full}} - P_{\mathrm{base}} = 0$ at every record:
$\max |\delta_0| = 6.39 \times 10^{-13}$ over the $900$ $N = 64$ records and
$8.94 \times 10^{-13}$ over the $40$ $N = 128$ records, so
$|\delta_0| < 10^{-12}$ at all $940$; the maximum over the full grid is
$8.94 \times 10^{-13}$, attained at $N = 128$. Per-$W$ values at the
primary point: $-3.77 \times 10^{-13}$ ($W = 2$), $+1.44 \times 10^{-13}$
($W = 3$), $0.0$ exactly ($W = 4, 6, 8$); the fit
$\delta_0(W) = \delta_\infty + c W^{-\kappa}$ degenerates to
$\delta_\infty = 0$ with zero residual, and $d\delta_0/d\Gamma = 0$. The optima
are not merely equal but achieved by identical optimal laws: the base-optimal law
itself satisfies the entire cubic block.

\paragraph{The exact $\varepsilon$ law.}
At asymptotic budgets the joint optimum obeys
$P(\varepsilon) = 5/6 - (2/3)\varepsilon$ --- the grid-decoupled corner of
Theorem~\ref{thm:corner} --- at every tested $\varepsilon$:

\begin{table}[h]
\centering
\begin{tabular}{@{}llll@{}}
\toprule
$\varepsilon$ & $P$ (asymptotic, LP) & $5/6 - (2/3)\varepsilon$ &
$|\delta_0|$ \\
\midrule
$0.02$ & $0.8200000000002683$ & $0.82$ & $2.8 \times 10^{-13}$ \\
$0.01$ & $0.8266666666669427$ & $0.8266666\ldots$ & $2.4 \times 10^{-13}$ \\
$0.005$ & $0.8300000000002841$ & $0.83$ & $2.2 \times 10^{-13}$ \\
$0.002$ & $0.8320000000003910$ & $0.832$ & $1.0 \times 10^{-13}$ \\
\bottomrule
\end{tabular}
\end{table}

Matched-variant values differ only through the finite-size budget center (e.g.\
$P = 0.8368627361$ at $\varepsilon = 0.002$, from
$B_1 = 84.712 < (4/3)\,64 = 85.33$). Moreover, adjoining the
$\lambda'$-Frobenius budget row alone to the two-moment baseline leaves the
optimum at $P_{\mathrm{base}}$ at every record: that row by itself adds
nothing.

\paragraph{Primary decision point}
(flat/flat, $\lambda' = 1/2$, $N = 64$, matched budgets, $\varepsilon = 0.05$,
$\Cled = 100$, coupled fuzz):
$P_{\mathrm{full}} = P_{\mathrm{base}} = 0.8050956815843511$,
$\delta_0 = -3.8 \times 10^{-13}$, identical optimal laws; the same $P$ at
fuzz $=$ none (the fuzz row is not load-bearing for the absorption). At
$N = 128$, $P_{\mathrm{full}} = P_{\mathrm{base}} = 0.8024488905983833$ (matched)
and $0.8000000000004773$ (asymptotic) at $\varepsilon = 0.05$, with
$\delta_0 = 0$ at all $40$ records.

\paragraph{The witness}
(\texttt{witness\_N64.json}; $N = 64$, matched budgets, $\varepsilon = 0.05$,
fuzz $=$ none). Three columns, all supported on the $65$-site $\psi_1$-zero grid
(spacing $64/65 = 0.9846153846$), marks $\{1, 2\}$ only, no pairs, $n_- = 0$:

\begin{table}[h]
\centering
\footnotesize
\begin{tabular}{@{}>{\raggedright\arraybackslash}p{3.1cm}
                  >{\raggedright\arraybackslash}p{3.0cm} l
                  >{\raggedright\arraybackslash}p{1.9cm}
                  >{\raggedright\arraybackslash}p{2.7cm} l@{}}
\toprule
weight & column & $\Nd$ & $F_1$ ($= \sum m^2$, exact) & $F'$ & $C'$ \\
\midrule
$0.5299906673474316$ &
vacancy lattice: $64$ simples on $64$ of the $65$ sites &
$64$ & $64$ & $125.0909091 = 4128/33$ & $245.4508724$ \\
\addlinespace
$0.16040139164388625$ &
$16$ doubles $+\ 32$ simples on an explicit $48$-site subset &
$48$ & $96$ & $153.7604346$ & $411.0865130$ \\
\addlinespace
$0.30960794100868216$ &
$32$ doubles on an explicit $32$-site subset &
$32$ & $128$ & $135.1959163$ & $301.3541285$ \\
\bottomrule
\end{tabular}
\end{table}

Row values in this table are the clean-integer-grid values
(\texttt{clean\_checks[].*\_clean} --- the canonical hand-checkable description);
the optimizer's stored law rows (\texttt{law[].F1/Fp/Cp}) differ from them by at
most $4.1 \times 10^{-5}$ per row (per column
$3.4 \times 10^{-6} / 1.6 \times 10^{-5} / 4.1 \times 10^{-5}$), and the
aggregates below are the stored law's, which saturate the budget edges exactly.
Aggregates: $E[F_1] = 88.94775275723633 = $ the upper Frobenius edge
$B_1(1 + \varepsilon)$ exactly; $E[F'] = 132.818$ (strictly inside its band
$[128.34, 141.85]$); $E[C'] = 289.32716513260516 = $ the lower cubic edge
$B_3(1 - \varepsilon)$ exactly (the doubles corner is cubic-short, per
$G(1/2) > 0$, and the arrangement supplies the budget from below);
$E[\Nd]/N = 51.52612362142/64 = 0.8050956815846875 \le 5/6 + 0.01$ (the
absorption bar of Section~\ref{sec:rowsystem}). Active rows at the optimum: only
the upper $F_1$ budget edge and the lower cubic edge; every ladder row slack by
orders of magnitude --- $E[n(2)] = 3.50$ and $E[n(3)] = 0.94$ against caps
$\Cled N V^{-4} = 400$ and $79.0$, all higher rows zero, and the vacancy column's
$n(V)$ vector identically zero. Every row was re-verified in exact
\texttt{Fraction} arithmetic (all $10$ checks pass, objective deviation
$5 \times 10^{-13}$), and all three columns were recomputed independently
through a position-space path (reproduction to $10^{-6}$ or better, feasibility
re-confirmed at $\Cled = 100$ and $\Cled = 10$).

The mechanism is explicit: $F_1 = \sum m^2$ exactly on all three columns
(Theorem~\ref{thm:parseval} --- the equality case of
Lemma~\ref{lem:diagpos}), so the bandwidth-one reading transfers the
\texttt{lemmaR\_tight} corner verbatim, while the sub-grid arrangement (which
sites, which clusters) freely meets the $\lambda'$ rows
(Proposition~\ref{prop:halfband}): the two doubles-subsets' occupancy-$2$
clustering supplies the cubic cross-terms ($C' = 411.09$ and $301.35$ against
isolated-block values $160$ and $256$) at zero cost in $\Nd$ and zero cost at
$\lambda = 1$. The witness uses no pairs, no marks $\ge 3$, no spectral escape ---
arrangement alone.

\paragraph{The role of pricing.}
No improving column (reduced cost $< -10^{-6}$) was found in $S = 200$ pricing
restarts at the primary and tightest-asymptotic points ($S = 100$ at $N = 128$,
recorded but with no stored verification block; $S = 40$ at $W$-scan
re-checks; $\min_{\mathrm{rc}} = 0.0$, $0$ improving columns at every stored
verification). The absorption conclusion does not rest on pricing optimality:
the witness is exhibited and verified, and
$P_{\mathrm{full}} = P_{\mathrm{base}}$ is an exact LP equality on the
accumulated dictionary in any case. Pricing convergence supports only the claim
that no configuration class outside the dictionary would carry a strictly
positive marginal value (and more columns can only lower both optima together,
the anti-absorption direction).

\paragraph{Robustness.}
The conclusion was stress-tested along eight independent axes: recomputation of
the witness; a per-column ladder at $\Cled = 4$, barely above the null floor
$3.82$; a missing-row hunt; near-CUE pinning; window-degeneracy probes; the
pair-channel attack; budget-center stress across $\varepsilon$ in
$0.002$--$0.10$ and $\Gamma$ to $0.25$; and seed and scale stability. Absorption
persists along all eight. Two of the axes produced further results --- on the
near-CUE class and the $p_1$ objective --- developed in
Sections~\ref{sec:nearcue} and~\ref{sec:p1}.

\section{The near-CUE class: absorption inside the unconditional data class}
\label{sec:nearcue}

The adversary of Section~\ref{sec:absorption} is free to hold
anti-correlated (non-CUE) bandwidth-one data. The strongest form of the no-go
pins the adversary inside the near-CUE class --- the data class of the
PairCeiling formalization accompanying~\cite{AF26} --- and asks whether the
cubic block has positive marginal value there. It does not, and the mechanism
is structural: at the pinned optimum the cubic row is slack
(Section~\ref{sec:slack}).

\subsection{The licensing theorem, and its three riders}
\label{sec:license}

The pinning rows $\bigl| E_w |c_j|^2 - j \bigr| \le \tau_2$ for
$1 \le j < N$ at $\lambda = 1$ (edge row free) are licensed by Theorem~1 of
Baluyot, Goldston, Suriajaya and Turnage-Butterbaugh~\cite{BGSTB24}
(arXiv:2306.04799 v1, June 2023; Acta Arithmetica, 2024), which is
unconditional and pointwise in $\alpha$, uniform on the closed band
$0 \le \alpha \le 1$:
\[
  F(\alpha) = T^{-2\alpha}(\log T + O(1)) + \alpha
    + O\bigl(1/\sqrt{\log T}\bigr),
\]
for the form factor defined with zeros counted with multiplicity and off-line
zeros entering with the depth weight $x^{\delta + \delta'}$ --- exactly the
$\cosh$ depth-weighting the model's form factor $c_j$ applies. With
$\alpha = j/N$, each open-band row is licensed pointwise per row,
unconditionally, in the depth-weighted multiplicity-counted form (directly, or
through their Lemma~5 admissible kernel class: even, $L^1$, supported in
$[-1, 1]$, Lipschitz at $0$ --- the bandwidth-one class); the diagonal spike is
confined to the edge rows the pinning class above leaves free, and the Poisson
factor acts below
the model's $1/N$ bin resolution whenever $N \ll \log T$. The uniformity is on
the closed band $0 \le \alpha \le 1$, not $|\alpha| < 1$: their Theorem~1
incorporates an earlier improvement, which they cite, extending the range ``up
to $\alpha = 1$ with explicit error terms''.

\paragraph{The three licensing riders (binding on every use of this section's
results):}

\begin{enumerate}
\item \textbf{Depth-aggregation only.} The licensed statistic never separates
  on-line from off-line mass: the Remark after their Theorem~2 states the method
  ``neither requires nor provides any information'' about whether the zeros are on
  the critical line. What is pinned is the $\cosh$-weighted aggregate $|c_j|^2$
  (pairs entering with depth weights), never an on-line-only statistic. Shedding
  the far-off-line coupling term costs an extra hypothesis --- the thin-box
  hypothesis or the strong zero-density hypothesis, which their Theorems~2--3 (the
  $61.7\%$ simple-zeros results) consume; the row class used here needs none of
  that.
\item \textbf{Finite-$T$ slack $O(N/\sqrt{\log T})$.} The licensed per-row slack
  at finite $T$ is $O(N/\sqrt{\log T})$ grid units (plus the edge spike), so
  $\tau_2 \in \{1, 4\}$ is the fixed-$N$, $T \to \infty$ asymptotic reading ---
  legitimate for the model's asymptotic semantics, but at finite $T$ the
  $O(1)$-unit pinning overstates the theorem until $\log T \gg N^2$.
\item \textbf{Second-moment scope.} The licensing theorem~\cite{BGSTB24} stops at
  the second moment --- pair
  correlation only; nothing in it licenses the cubic budget column, whose
  $c_3(\lambda')$ identification remains open (the cubic budget center,
  Section~\ref{sec:flags}).
\end{enumerate}

In sum: the pinning rows are an unconditional asymptotic data class, pointwise
on the open band in the depth-weighted form, and the \emph{strength} of the
near-CUE conclusion --- absorption inside the two-moment certificate's own data
class --- inherits riders 1--3.

\subsection[delta\_0' = 0 exactly, both N, both centerings]{$\delta_0' = 0$
exactly, both $N$, both centerings}
\label{sec:t2result}

Row class: all the rows of Section~\ref{sec:rowsystem} (fuzz $=$ none --- the
harsher variant for absorption)
plus the pinning rows at $\tau_2$. Four cells at $N = 64$ (matched budgets,
$\Cled = 100$, $W \le 8$):

\begin{table}[h]
\centering
\small
\begin{tabular}{@{}llllr >{\raggedright\arraybackslash}p{4.3cm}@{}}
\toprule
$\tau_2$ & $\varepsilon$ & $P_{\mathrm{full}}^{\mathrm{pin}}$ &
$P_{\mathrm{base}}^{\mathrm{pin}}$ & $\delta_0'$ & reading \\
\midrule
$1$ & $0.05$ & $0.8332334059839699$ & $0.8332334059839699$ &
$\mathbf{0.0}$ (exact) & absorption (exact) \\
\addlinespace
$1$ & $0.02$ & $0.8359542395410772$ & $0.8359374423865479$ &
$1.68 \times 10^{-5}$ &
absorb ($0$ within the residual band
$[-3.1 \times 10^{-5}, +1.8 \times 10^{-5}]$) \\
\addlinespace
$4$ & $0.05$ & $0.8150864520465212$ & $0.8137831067863290$ &
$1.30 \times 10^{-3}$ & bounded residual, unconverged upper bound \\
\addlinespace
$4$ & $0.02$ & $0.8316328692004420$ & $0.8301979402692543$ &
$1.43 \times 10^{-3}$ & bounded residual, unconverged upper bound \\
\bottomrule
\end{tabular}
\end{table}

At the primary cell ($\tau_2 = 1$,
$\varepsilon = 0.05$), $\delta_0' = 0$ is an exact LP equality --- identical
optima --- at every one of $9$ cumulative $S = 200$ verification passes across a
${\sim}1{,}100$-column dictionary enrichment; the common level fell
$0.8333622 \to 0.8332334$ while the equality never broke (residual improvements
are common-mode), and the reduced-cost bound for the true marginal value is
$|\delta_0'| \le 1.2 \times 10^{-5}$. Stability under budget re-centering: the
asymptotic variant gives
$P_{\mathrm{full}} = P_{\mathrm{base}} = 0.8332332347490976$,
$\delta_0' = 0.0$ exact.

\paragraph{$N = 128$ confirmation.}
The pinned LP was infeasible on the raw dictionary and was bootstrapped with
$300$ CUE and doubles-decorated-CUE columns (anti-absorption-safe), then built
out over $10$ cumulative $S = 200$ passes to $6{,}457$ columns at the matched
cell ($6{,}739$ after the asymptotic re-centering):
$P_{\mathrm{full}} = P_{\mathrm{base}} = 0.8374218534574983$ with
$\delta_0' = 0.0$ exact at the matched cell, and $0.8373740026259704$ with
$\delta_0' = 0.0$ exact at the asymptotic re-centering. Build-out invariance:
$\delta_0' \le 8 \times 10^{-16}$ at every recorded dictionary re-solve
($1{,}842 / 6{,}457 / 6{,}739$ columns) --- the conclusion is independent of the
convergence level. Level trend: $P^{\mathrm{pin}} = 0.8374$ ($N = 128$) vs
$0.8332$ ($N = 64$), both ${\sim}5/6 - O(10^{-3})$, the pinned-baseline
finite-size trend.

\subsection{Absorption by slack}
\label{sec:slack}

The near-CUE witnesses explain the mechanism. At $N = 64$: a $64$-column
law, marks $\{1, 2\}$ only, no pairs, no marks $\ge 3$, with
$E|c_j|^2 = j + 1$ for every $j = 1, \ldots, 63$ to $2 \times 10^{-12}$ (maximum
deviation $1.84 \times 10^{-12}$, at $j = 62$) --- all $63$ upper pinning rows
active at exactly the $+\tau_2$ allowance --- and the only active rows at
the pinned optimum are the $63$ pinning rows (and mass): $F_1$, $F'$, $C'$, and
every ladder row are strictly interior. The cubic block does not bind.
The $+\tau_2$ allowance pays for the doubles; the cubic budget is met with room
to spare. \texttt{Fraction} re-verification passes on the full row class (worst
pinning margin $-1.0 \times 10^{-9}$, inside the documented $10^{-6}$ grace). At
$N = 128$ the same structure: $128$ columns, marks $\{1, 2\}$, no pairs, upper
pinning edge ridden at $\tau_2 - 10^{-5}$ (an interior shrink absorbing a
solver-tolerance overshoot; the \texttt{Fraction} re-verification against the
true $\tau_2 = 1$ box passes with margin $+1.0 \times 10^{-5}$ strictly
interior), $E[n_{\mathrm{simple}}]/N = 0.6748$.

This is a stronger statement than absorption in the unpinned class: inside the
near-CUE class
the cubic row is not absorbed by an adversary straining against it --- it is
slack at the optimum. There is no escape route inside the certificate's own data
class, and the cubic block has no positive marginal value there, at both lattice
sizes, at both budget centerings, subject to riders 1--3 of
Section~\ref{sec:license}.

\subsection[The tau\_2 = 4 annotation]{The $\tau_2 = 4$ annotation}

At loose pinning ($\tau_2 = 4$) the cubic row is active and the residuals
$\delta_0' = 1.30 \times 10^{-3} / 1.43 \times 10^{-3}$
($\varepsilon = 0.05 / 0.02$) exceed their pricing bands
($[0.9 \times 10^{-3}, 1.5 \times 10^{-3}]$) and were still declining under
enrichment ($1.28 \times 10^{-3} \to 1.17 \times 10^{-3}$ over five passes): they
are unconverged upper bounds at loose pinning, an order of magnitude below the
$10^{-2}$--$3 \times 10^{-2}$ $\delta_0$ the proposal projected, obtained from
shortened runs and reported as such. They are diagnostic annotations, not
certificate claims, and they sharpen the scoping: whatever residual cubic signal
exists at this scale lives
strictly between the near-CUE class ($\tau_2 = 1$: exact zero by slack) and the
free unpinned class (exact zero by witness), i.e.\ in loosely pinned intermediate
worlds, at magnitudes far below the scale at which the question was posed.

\section{The simple-fraction benchmark}
\label{sec:p1}

The headline of the two-moment certificate~\cite{AF26} is the simple-zero
fraction ($2/3$ at flat
window; $0.6725$ window-optimized; ceiling $0.6818287$), so a no-go scoped only
to $\Nd$ would be incomplete. The
computation was therefore re-run with the objective $\min E_w[p_1]$,
$p_1(c) = (\#\text{mark-}1\ \mathrm{atoms}
 + 2 \times \#\text{mult-}1\ \mathrm{pairs})/N$, on the dictionary as enriched by
the near-CUE runs of Section~\ref{sec:nearcue} ($5{,}134$ columns), with the
stopping criterion --- no improving column in $S = 200$ pricing restarts ---
applied to both systems:

\begin{table}[h]
\centering
\footnotesize
\begin{tabular}{@{}>{\raggedright\arraybackslash}p{1.7cm}
                  >{\raggedright\arraybackslash}p{2.9cm}
                  >{\raggedright\arraybackslash}p{2.9cm}
                  >{\raggedright\arraybackslash}p{1.4cm}
                  >{\raggedright\arraybackslash}p{3.2cm}
                  >{\raggedright\arraybackslash}p{2.2cm}@{}}
\toprule
cell & $p_{1,\mathrm{full}}$ & $p_{1,\mathrm{base}}$ & $\delta_{p_1}$ &
analytic corner $2 - (B_1/N)(1 + \varepsilon)$ & converged \\
\midrule
matched, $\varepsilon = 0.05$ & $0.6101913631697143$ &
$0.6101913631684138$ & $1.3 \times 10^{-12}$ & $0.6101913631681821$ &
yes: $S = 200$ clean, $\min_{\mathrm{rc}} = 0.0$, both systems \\
\addlinespace
asymptotic, $\varepsilon = 0.002$ & $0.6640000000012899$ &
$0.6640000000004871$ & $8.0 \times 10^{-13}$ &
$0.6640000000000001 = 2 - (4/3)(1.002)$ &
yes: $S = 200$ clean, both systems \\
\bottomrule
\end{tabular}
\end{table}

These are the only cells reported here in which the pricing stopped clean for
both systems at once: zero improving columns at $S = 200$ for the full and base
systems. Cross-checks on the final dictionary:
coupled fuzz at the primary cell $\delta_{p_1} = 2.9 \times 10^{-13}$; $W$-scan
over $\{2, 3, 4, 6, 8\}$ max $|\delta_{p_1}| = 1.3 \times 10^{-12}$ (at $W = 8$;
all $< 1.4 \times 10^{-12}$); the labeled on-line-only diagnostic variant
$3.3 \times 10^{-13}$. The optimum is the exact marks-$\{1,2\}$ doubles corner of
Theorem~\ref{thm:corner}(b) --- pairs cannot help (a mult-$1$ pair buys $2$
simple points at $\ge$ twice the two-simples Frobenius cost; a mult-$2$ pair is
dominated by two doubles) --- and the $\varepsilon \to 0$ value is
$2 - 4/3 = \mathbf{2/3}$, the model's bandwidth-one analog of Montgomery's
simple-zeros corner. The witness laws are the same $\psi_1$-grid family as the
witness of Section~\ref{sec:certification}, reweighted per cell, rational
verification passing.

Two scoping notes. (i) The cited $0.6818287$ ceiling belongs to the
\emph{window-optimized} certificate class (the Theorem-D effect), which the
flat/flat row system used here does not consume; the flat-window model corner is
$2/3$, and what is proved here is that the cubic block does not move the
flat/flat optimum for either objective, leaving the $0.6818287$ ceiling
untouched. (ii) With this section, the no-go covers both benchmarks: the
distinct-zeros corner $5/6$ and the simple-fraction corner $2/3$, each with
$\delta = 0$ at convergence.

\section{Capacity and vacuity: no spectral escape}
\label{sec:capacity}

Two theorems close the spectral-escape routes from opposite ends, and together
with the decision program they complete the pricing: a scalar cubic tail row
with any
divergent cutoff is unconditionally vacuous (Theorem~\ref{thm:garnish} --- the
half of the no-go that needs no linear program at all); and under the all-$V$
ladder the same channel is capped at
$o(N)$ with a sharp constant (Theorem~\ref{thm:capacity}). What remains is
position/interference freedom --- which is exactly the channel the witness of
Section~\ref{sec:certification} uses.

\subsection{The sharp capacity theorem}
\label{sec:sharpcap}

\paragraph{Setting (the continuum garnish LP).}
Heights normalized so the ladder's absolute floor is $1$. Spectral escape is a
positive measure $\mu$ on $[1, \infty)$ --- $\mu([h, \infty))$ is the count per
$N$ of escaped eigenvalue mass at height $\ge h$ --- subject to
\begin{align*}
  \text{(Frobenius slack)} \quad & \int h^2\, d\mu(h) \le \varepsilon_g \,,\\
  \text{(all-$V$ ladder)} \quad & \mu([h, \infty)) \le C h^{-4}
    \quad \text{for every } h \ge 1 \qquad (C = \Cled),
\end{align*}
and the objective is the absorbable cubic charge $\int h^3\, d\mu(h)$.
(Equivalently, in tail-function form $N(h) = \mu([h, \infty))$: maximize
$N(1) + 3\int_1^\infty h^2 N\, dh$ subject to
$N(1) + 2\int_1^\infty h N\, dh \le \varepsilon_g$ and $N \le C h^{-4}$; the
forms are identical by integration by parts.)

\begin{theorem}[sharp capacity]
\label{thm:capacity}
For every $C > 0$ and $0 < \varepsilon_g \le 2C$:
\[
  \max_\mu \int h^3\, d\mu = 2\sqrt{2}\,\sqrt{C \varepsilon_g}\,,
\]
attained by the min-profile $N_*(h) = C\min(a^{-4}, h^{-4})$,
$a = \sqrt{2C/\varepsilon_g}$ (as a measure: no mass on $[1, a)$; density
$4 C h^{-5}\, dh$ on $(a, \infty)$).
\end{theorem}

\begin{proof}
Upper bound: for any feasible $\mu$ and any $a > 0$, the pointwise inequality
$h^3 \le a h^2 + (h^3 - a h^2)_+$ on $[1, \infty)$ gives
\[
  \int h^3\, d\mu \le a \varepsilon_g + \int (h^3 - a h^2)_+\, d\mu \,.
\]
Writing $g(h) = (h^3 - a h^2)_+ = \int_a^h g'(s)\, ds$ for $h \ge a$, with
$g'(s) = 3 s^2 - 2 a s > 0$ on $(a, \infty)$, Tonelli and the ladder give
\[
  \int g\, d\mu = \int_a^\infty g'(s)\, \mu([s, \infty))\, ds
  \le C \int_a^\infty (3 s^2 - 2 a s) s^{-4}\, ds
  = C\Bigl( \frac{3}{a} - \frac{a}{a^2} \Bigr)
  = \frac{2C}{a} \,.
\]
So $\int h^3\, d\mu \le a \varepsilon_g + 2C/a$ for every $a > 0$; minimizing
over $a$ (minimum at $a = \sqrt{2C/\varepsilon_g}$) gives
$2\sqrt{2 C \varepsilon_g}$. Attainment: $\varepsilon_g \le 2C$ makes
$a \ge 1$, so $N_*$ is an admissible tail function; its Frobenius charge is
$C a^{-4} + C a^{-4}(a^2 - 1) + C a^{-2} = 2C/a^2 = \varepsilon_g$ exactly, and
its cubic charge is $C/a + 3C/a = 4C/a = 2\sqrt{2}\sqrt{C \varepsilon_g}$,
meeting the bound.
\end{proof}

\begin{remk}[the two-sided squeeze, made exact]
The optimal dual split is
\[
  [a \times \text{Frobenius}] + [\text{ladder integral above } a] :
\]
bounded heights are blocked by the Frobenius cost (mass at height $h$ buys $h^3$
of cubic but pays $h^2$ of Frobenius, ratio $h \le a$), divergent heights by the
ladder cap. The absolute floor $C_{\mathrm{abs}}$ (normalized to $1$) appears
nowhere in the upper bound --- only in the attainability condition $a \ge 1$ ---
which is why a criterion of the form ``the marginal value is positive iff
$2\Cled/C_{\mathrm{abs}} < 4/3$'' cannot decide the question: sharp constants at
the floor are irrelevant against this adversary class.
\end{remk}

\begin{proposition}[the naive profile is infeasible]
\label{prop:naive}
The point-mass-plus-tail profile (a ladder-saturating atom of size
$n_0 = C b^{-4}$ at height $b = \sqrt{3C/\varepsilon_g}$, plus the
ladder-saturating tail above $b$) has Frobenius charge exactly $\varepsilon_g$
and cubic charge
$5C/b = (5/\sqrt{3})\sqrt{C \varepsilon_g} = 2.8868\sqrt{C \varepsilon_g}$ ---
but it violates the cumulative ladder on the open interval $(2^{-1/4} b, b)$,
where its count $2 C b^{-4}$ exceeds the cap $C h^{-4}$. The true optimum is
Theorem~\ref{thm:capacity}'s $2\sqrt{2} = 2.8284271247\ldots < 5/\sqrt{3}$; the
earlier estimates $5/\sqrt{3}$ and the two-term dyadic $2\sqrt{C\varepsilon}$ are
of the same order and are superseded. The constant is confirmed three independent
ways: the min-profile derivation above, an analytic re-derivation together with
a linear program ($2.8244$ at $h_{\max} = 400$, gap $=$ the predicted truncation
tail $3/h_{\max}$), and a separate numerical optimum.
\end{proposition}

\begin{proof}
The tail part has density $4 C h^{-5}$ on $(b, \infty)$. Frobenius: atom
$C b^{-2}$ $+$ tail $2 C b^{-2} = 3 C b^{-2} = \varepsilon_g$ at the stated $b$.
Cubic: atom $C/b$ $+$ tail $4C/b$. Ladder violation: for
$2^{-1/4} b < h < b$ the cumulative count is $n_0 + C b^{-4} = 2 C b^{-4}$ while
the cap is $C h^{-4} < 2 C b^{-4}$ throughout the open interval.
\end{proof}

\begin{corollary}[capacity is $o(N)$: the positive complement of the no-go]
\label{cor:oN}
In the model's normalized units, the maximal cubic shift purchasable by
spectral escape under
the row system with the all-$V$ ladder in force is
\[
  |\Delta \text{cubic}| \le 2\sqrt{2}\,\sqrt{\Cled \varepsilon_g}\, N = o(N)
  \qquad \text{as } \varepsilon_g \to 0 \text{ at fixed } \Cled
\]
(e.g.\ $\Cled = 1000$, $\varepsilon_g = 10^{-8}$: capacity $0.0089 N$). A
doubles-vs-pairs-scale swing ($\Theta(N)$, e.g.\ $(4/3)N$) can never be produced
by spectral escape once the all-$V$ ladder is consumed: any absorption of the
cubic block must run through position/interference freedom --- the channel
Theorem~\ref{thm:parseval} exhibits and the witness uses. (In the decision LP this
theorem enters as the coupled fuzz rows, tangent-cut on the concave $\sqrt{\ }$;
the fuzz row is not binding at the optimum.)
\end{corollary}

Numerical confirmation (\texttt{verify\_t5.py}, with results in
\texttt{verify\_t5\_}\allowbreak\texttt{out.json}): attainment profile exact
to $4 \times 10^{-11} / 9 \times 10^{-13}$; an independent discretized LP over
point masses converges to $0.4000$ from below under joint grid-and-truncation
refinement --- tail-corrected values $0.3950 / 0.3985 / 0.3991$ at
($600$ points, $h_{\max} = 2000$), ($2000$ points, $h_{\max} = 2000$),
($4000$ points, $h_{\max} = 8000$), $C = 1$, $\varepsilon_g = 0.02$, each the raw
LP value plus the predicted truncation tail $3C/h_{\max}$ --- never exceeding the
bound; the naive profile's value $0.408248$ and its ladder violation are
reproduced; duality and scaling verified across $(C, \varepsilon)$ over five
orders of magnitude.

\subsection{Divergent-cutoff vacuity: the garnish theorem}

\paragraph{Setting (the row system in the divergent-cutoff form the proposal
stated; family semantics).}
$N = N(T) \to \infty$; $V_0 = V_0(T) \to \infty$ any divergent cutoff with
$V_0 = o(N^{1/3})$ (any polylog cutoff qualifies --- in particular the proposal's
$(\log\log T)^3$ and every $(\log T)^A$). The as-written system consumed its
equalities ``with $o(N)$ precision'' --- no fixed tolerance rate --- so
feasibility is a property of families $(c_T)$ of model configurations (no mark cap
was imposed as written), with spectral data read from isolated-block spectra,
admissible iff
\begin{itemize}
\item[(W1)] mass: $\sum m = N$ (exact);
\item[(W2)] Frobenius: $F = \kappa N + o(N)$, $\kappa = \kappa(1) = 4/3$ (the
  as-written system's bandwidth-one Frobenius row);
\item[(W3)] signed cubic: $C = c_3 N + o(N)$;
\item[(W4)] scalar tail row:
  $\sum_{|\lambda_i| \ge V_0} |\lambda_i|^3 = o(N)$;
\item[(W5)] count/inertia rows imposed only at thresholds $V \ge V_0$, plus
  $n_-$ at most the number of pairs;
\end{itemize}
objective $\liminf \Nd/N$ (or $\liminf p_1$). (A fixed tolerance rate shrinking
faster than $1/V_0$ would be a stronger system than the one written: the
consumed prime-side equalities carry no such rate, so imposing one would assume
a stronger input than the sources provide.)

\begin{theorem}[garnish absorption]
\label{thm:garnish}
Let $(c_T)$ satisfy (W1), (W2), (W4), (W5), contain at least $s_0 N$ simple
on-line atoms for a fixed $s_0 > 0$ (a simple-atom fraction, unrelated to the
marginal value $\delta_0$), and fall short of the (W3) cubic equality by
$\Delta_T N$ with $0 \le \Delta_T \le \Delta_{\max}$ fixed. Put
$h_0 = \lfloor V_0/2 \rfloor$, $n_0 = \lceil \Delta_T N/h_0^3 \rceil$, and let
$c'_T$ be $c_T$ with $n_0$ additional isolated on-line atoms of mark $h_0$ (at
unoccupied, well-separated positions) and $n_0 h_0$ simple atoms deleted. Then
for all large $T$, $c'_T$ satisfies every row of the as-written system, and
\[
  \Nd(c'_T) = \Nd(c_T) - n_0 (h_0 - 1) \le \Nd(c_T),
  \qquad p_1(c'_T) \le p_1(c_T).
\]
Row-by-row (per $N$; every shift $o(1)$): mass exact $0$; Frobenius
$+\ \Delta_T/h_0 + O(1/h_0^2 + h_0^2/N) = o(1)$; cubic $+\ \Delta_T + o(1)$
(landing the (W3) equality); tail row $+\ 0$ (the garnish sits at $h_0 < V_0$);
count rows at $V \ge V_0$: $+\ 0$; inertia unchanged (only positive on-line
eigenvalues added).
\end{theorem}

\begin{proof}
Feasibility: $n_0 h_0 \le \Delta_{\max} N/h_0^2 + h_0 = o(N) < s_0 N$ for large
$T$ --- enough simples exist to delete. Each garnish atom is an isolated on-line
atom of mark $h_0$: eigenvalue $h_0$, Frobenius $h_0^2$, cubic $h_0^3$, mass
$h_0$. Mass: adds and deletes $n_0 h_0$ exactly. Cubic: adds $n_0 h_0^3$ in
$[\Delta_T N, \Delta_T N + h_0^3]$, subtracts
$n_0 h_0 \le \Delta_{\max} N/h_0^2 + h_0$; net
$\Delta_T N + O(h_0^3) + O(N/h_0^2) = \Delta_T N + o(N)$ since
$h_0 \to \infty$ and $h_0^3 = o(N)$ --- (W3) now holds. Frobenius: adds
$n_0 h_0^2 \le \Delta_{\max} N/h_0 + h_0^2 = o(N)$, inside (W2)'s window. Tail
and count rows: every garnish eigenvalue equals $h_0 = \lfloor V_0/2 \rfloor <
V_0$, and all thresholds sit at $V \ge V_0 > h_0$, so the added spectrum is
invisible; deletions only decrease those rows. Objectives: $\Nd$ changes by
$+\,n_0 - n_0 h_0 = -n_0(h_0 - 1) \le 0$; the deleted atoms were simple and the
added ones have mark $h_0 \ge 2$, so $p_1$ decreases.
\end{proof}

\begin{remk}[the growth cap is removable]
The hypothesis $V_0 = o(N^{1/3})$ is a convenience, not a constraint: for an
arbitrary divergent cutoff, cap the garnish height at
$h_0 = \min(\lfloor V_0/2 \rfloor, \lfloor N^{1/4} \rfloor)$. Then
$h_0 \to \infty$, $h_0^3 \le N^{3/4} = o(N)$, and $h_0 < V_0$ still holds, so
every estimate in the proof goes through verbatim --- the vacuity covers any
divergent cutoff.
\end{remk}

\begin{corollary}[vacuity of the scalar divergent-cutoff cubic row]
\label{cor:vacuity}
Suppose the two-moment system $\{$(W1), (W2)$\}$ admits optimizer families with a
fixed positive simple fraction, spectra bounded by a fixed constant, and bounded
cubic deficit ($C(c_T) \le c_3 N + o(N)$, deficit $\le \Delta_{\max} N$). Then,
for both benchmarks, the infimum over the as-written system equals the infimum
over the two-moment system: the marginal value of the entire block
$\{$(W3) $+$ (W4) $+$ (W5)$\}$ is zero --- $\delta_0 = 0$ as written. The deficit
hypothesis is the actual situation: the corner families of
Theorem~\ref{thm:corner} have marks in $\{1, 2\}$ (spectra bounded by $2$, simple
fraction $2/3$ in the limit) and undershoot the cubic budget
(Corollary~\ref{cor:anchors}: $G(1/2) = +1/2 > 0$ --- the null budget sits
strictly above the doubles plane, so doubles-saturated configurations are
cubic-short; the witness accordingly sits at the lower cubic edge).
Whatever cubic demand the row was hoped to impose on the two-moment-degenerate
corner --- up to any fixed $\Delta_{\max}$, including the full ``$(4/3)N$
doubles-vs-pairs swing'' of the proposal --- the garnish supplies it invisibly,
at an objective decrease.
\end{corollary}

\begin{proof}
One direction is trivial (more rows can only raise the infimum). Conversely, take
the corner families: bounded spectra plus $V_0 \to \infty$ make (W4) and the
$V \ge V_0$ count rows identically zero for large $T$, so they satisfy (W1),
(W2), (W4), (W5); apply Theorem~\ref{thm:garnish} to land (W3); the garnished
family is admissible for the whole system at a weakly smaller objective.
\end{proof}

\begin{remark}[what this does and does not say]
\label{rem:scope}
\begin{enumerate}
\item \textbf{It kills the scalar row, not the underlying idea.} The vacuity is
  about the as-written system, whose count/tail information starts only at
  $V_0$. The proposal's own Chebyshev proof in fact yields the all-$V$ ladder
  $n(V) \le \Cled N V^{-4}$ for every $V \ge C$ at $\vartheta < 1$ --- more than
  the divergent-cutoff form in which it was stated; under that system the
  garnish channel is not free but capped, and the cap is exactly
  Theorem~\ref{thm:capacity}'s $2\sqrt{2}\sqrt{\Cled \varepsilon_g}\, N = o(N)$.
  Claiming exactly what that proof yields is precisely the difference between
  Corollary~\ref{cor:vacuity} and Theorem~\ref{thm:capacity}.
\item \textbf{The garnish needs unbounded marks}
  ($h_0 = \lfloor V_0/2 \rfloor \to \infty$): it lives outside every fixed-$W$
  alphabet --- which is why the adversary class here carries a divergent mark
  alphabet $W$, and why a violation found only at marks $\{1, 2\}$ would not
  settle the question. (A pair-based
  garnish with slowly divergent pair multiplicity reaches the same conclusion
  through Proposition~\ref{prop:pairblock}'s deep-pair blocks.)
\item \textbf{The three-part structure of the no-go.}
  Theorem~\ref{thm:garnish}/Corollary~\ref{cor:vacuity} dispose of spectral
  escape under the divergent-cutoff row; Theorem~\ref{thm:capacity} caps it
  under the all-$V$ ladder; and the
  computational content (Sections~\ref{sec:absorption}--\ref{sec:p1}) is that
  the remaining channel --- position/interference freedom --- absorbs the cubic
  block too, at $\delta_0 = 0$ exactly. Together these are the sharpened no-go.
\end{enumerate}
\end{remark}

Numerical confirmation (\texttt{verify\_t6.py}): the shift table at
$\Delta = 4/3$, $V_0 = (\log\log T)^3$ over a $\log T$ grid
extended to $10^6$ --- Frobenius shift $4/(3 h_0) = 0.0992$ down to
$0.0010 \to 0$, $\Nd$ shift negative throughout, cubic shift $+4/3$ exactly, tail
and count rows $0$ --- reproducing two independent computations of the same
table; an abstract LP demonstration in which the cubic equality is infeasible
without the garnish variable (which would otherwise read as a positive marginal
value) and returns to the two-moment corner $5/6 - O(\mathrm{tol})$ with it; and
an exact-rational instance with all shifts exact.

\section{Scope, residuals, and what survives}
\label{sec:scope}

This section states where $\delta_0 = 0$ is exact, where it is a bounded
residual, what is finite-scale, and what remains open, and closes with the one
range of configurations in which the cubic row does carry a positive marginal
value.

\subsection{The residual table}

The exact zero is flat/flat-specific. Theorem~\ref{thm:parseval}'s
mechanism says exactly why: the flat bandwidth-one kernel's zeros are equally
spaced, so one site grid kills all $\lambda = 1$ cross-terms simultaneously; a
cosine window's kernel zeros are not equally spaced, so no grid does, and the
degeneracy breaks by a microscopic but nonzero amount. The complete residual map:

\begin{table}[h]
\centering
\small
\begin{tabular}{@{}>{\raggedright\arraybackslash}p{4.0cm}
                  >{\raggedright\arraybackslash}p{3.7cm}
                  >{\raggedright\arraybackslash}p{6.3cm}@{}}
\toprule
configuration & residual $\delta_0$ & status \\
\midrule
flat/flat, entire $940$-record decision grid & $0$ ($< 10^{-12}$) &
exact; witness-backed \\
\addlinespace
flat at $\lambda = 1$ / $\cos(1.6 s)$ at $\lambda'$ &
$5.95 \times 10^{-14}$ ($\varepsilon\ .05$),
$-1.14 \times 10^{-13}$ ($\varepsilon\ .02$) &
exact-zero class: changing the $\lambda'$ window moves nothing --- the
absorber's freedom is positional, not spectral \\
\addlinespace
$\cos(\sqrt{2}\, s)$ [Montgomery--Taylor] at $\lambda = 1$ / flat &
$+6.41 \times 10^{-5}$ ($\varepsilon\ .05$),
$+2.28 \times 10^{-5}$ ($\varepsilon\ .02$) &
bounded residual: the MT cosine breaks the exact grid degeneracy; unconverged
upper bounds, three orders below the $\varepsilon$-slack scale \\
\addlinespace
$\cos(\sqrt{2}\, s)$ [MT] / $\cos(1.6 s)$ &
$+6.27 \times 10^{-5}$ ($\varepsilon\ .05$),
$+1.91 \times 10^{-5}$ ($\varepsilon\ .02$) & same class \\
\addlinespace
near-CUE $\tau_2 = 1$ (certified, $N = 64$ and $128$, both centerings) &
$0$ exact (reduced-cost bound $\le 1.2 \times 10^{-5}$ at $N = 64$) &
absorption by slack (Section~\ref{sec:slack}) \\
\addlinespace
near-CUE $\tau_2 = 4$ (loose box, diagnostic) &
$1.30 \times 10^{-3} / 1.43 \times 10^{-3}$, pricing bands
$[0.9 \times 10^{-3}, 1.5 \times 10^{-3}]$ &
unconverged upper bounds, an order below the $10^{-2}$--$3 \times 10^{-2}$
scale at which the question was posed; the cubic row is active here \\
\addlinespace
independent tight-$\varepsilon$ near-CUE probes (dictionary-limited, no column
generation) & $\le 5 \times 10^{-4}$, non-monotone in $\varepsilon$ &
upper bounds, same order as the documented ${\sim}3 \times 10^{-4}$ pricing
uncertainty --- the regime where any residual signal would concentrate \\
\bottomrule
\end{tabular}
\end{table}

Every nonzero entry is an unconverged upper bound on the converged marginal
value, and every one sits at least an order of magnitude below the scale at
which the question was posed. The scope of the exact-zero claim is therefore
this: $\delta_0 = 0$
exactly at the primary flat/flat configuration; $\le O(10^{-4})$ across the
window family and $\le O(5 \times 10^{-4})$ under near-CUE pinning at that scale,
within pricing-convergence uncertainty; and at $\tau_2 = 1$ the near-CUE bound
is replaced by the exact zeros of Section~\ref{sec:nearcue}.

\subsection{Finite scale}

The computational certification is at $N = 64$ and $N = 128$, with the same
conclusion at both (and across two budget centerings whose cubic centers differ
by $5\%$ --- the $F_1$ centers by $0.7\%$, $F'$ by $2.6\%$). It is not
extrapolated beyond that: the continuum analog of the decoupling mechanism is
exact (Proposition~\ref{prop:lattice} --- the integer lattice under the sinc
kernel, with the finite-$N$ grid its exact periodization), so the absorption has
a structural reason to persist at every scale, and the model-level theorems of
Sections~\ref{sec:structure},
\ref{sec:atomonly}--\ref{sec:pairchannel}, and~\ref{sec:capacity} are
$N$-generic (the pair-channel capacity constants were certified at $N = 64$, the
primary geometry; the $N = 128$ re-run
reproduces every identity and inequality --- see (O4) below). No claim about $N$
beyond $128$ rests on computation.

\subsection{Open questions and residual coverage}
\label{sec:flags}

\begin{itemize}
\item \textbf{The cubic budget center.} The identification
  $c_3(\lambda') = m_3(\lambda')$ ($= 5$ at $\lambda' = 1/2$) rests on
  Rudnick--Sarnak-range GUE $3$-level correlations and on the polylog-window
  uniformity of the diagonal method of~\cite[\S5]{AF26};
  Theorem~\ref{thm:closedforms} proves the sine value, not the identification.
  The absorption conclusion is insensitive to it: $\delta_0 = 0$ across
  $\varepsilon$ in $0.002$--$0.10$ and $\Gamma$ to $0.25$, and across budget
  centers differing by $5\%$ --- an $O(1)$-per-zero error in $c_3$ cannot flip it
  (one of the robustness checks recorded in Section~\ref{sec:certification}).
  Nothing in~\cite{BGSTB24} licenses anything at third order (rider~3).
\item \textbf{The ladder's provenance.} The all-$V$ ladder consumed by
  Theorem~\ref{thm:capacity}'s system rests on the all-$V$ spectral-tail
  statement, whose step from a bound above one fixed cutoff to a bound valid at
  every threshold $V$ is not proved. Assuming the ladder only
  strengthens the no-go (absorption holds even with
  the ladder in force), and Theorem~\ref{thm:garnish} needs no ladder at all;
  nothing in this paper depends on that bridge.
\item \textbf{Heuristic pricing.} For classes without an exhibited witness (the
  nonzero rows of the residual table; the claim that no undiscovered
  configuration class carries a strictly positive marginal value), the evidence
  is converged column-generation pricing
  ($S = 200$ restart protocol), which is heuristic adversary-optimality; it is
  never load-bearing for the exact-zero claims, which are witness-backed LP
  equalities.
\item \textbf{Pair-channel residue} (Section~\ref{sec:pairchannel} and the
  companion note cited there): (O1) The crowding cap's
  ledger constant is interval-certified to $w \le 0.98$ ($0.98465 < 1$); the
  full-family ratio $0.9775$ is a grid artifact --- on
  $w \in (0.98, 1]$ the ledger fails (certified $> 1$) and coverage is
  adversarial numerics ($+17.46$, re-run at the sliver) plus the $8/9$ backstop;
  the cap itself on $w \in (0.82, 0.98]$ remains granted-not-proved (its
  $w \le 0.82$ self-consistency is interval-certified on both sides). (O2)
  Single-pair depths in $(0.45, 1.0)$ at unrestricted mass, multi-pair columns
  at the deep-probe depths $w \in \{1.5, 2\}$, and multi-pair configurations
  outside the admissible class (depth beyond $w = 2$, or unbounded mass), are
  covered by numerics plus the $8/9$ backstop only --- not needed for any LP
  statement. (O3) The sharp conjecture
  (Section~\ref{sec:pairchannel}) is open; near-tight configurations ($+0.03$)
  show that the constant is not easily improved. (O4) The certified capacity
  constants are at $N = 64$; the $N = 128$ re-run
  reproduces every identity and inequality with $\le {\sim}2.5\%$ drift in the
  headline certified constants (deep-tail table values, $y \ge 1.5$, drift up to
  ${\sim}8\%$), moves the deep-pair threshold of Theorem~\ref{thm:singlepair} from
  $d \ge 1.0$ to $d \ge 1.1$, and shows the
  near-endpoint ledger failure at $N = 128$ as well (smaller sliver, better capped
  constant $0.9633$). (O5) All pair-channel constants are for the flat
  $\lambda = 1$ window; none transfer to the MT-cosine window without
  recomputation --- consistent with that window's residual row above.
\item \textbf{The formalized ceilings.} The $0.6818287$
  bandwidth-one ceiling belongs to the window-optimized certificate class; the
  $0.8453$ trace/Frobenius cap stands on $\kappa \ge 2/\sqrt{3}$. The flat/flat
  row system used here consumes neither, and neither is revisited here: the
  claim made here is that cubic augmentation moves
  nothing at the flat/flat operating point and produces only the bounded residuals
  above elsewhere.
\end{itemize}

\subsection[Outlook: the lambda' payoff curve (outside the proven
ladder)]{Outlook: the $\lambda'$ payoff curve (outside the proven
ladder)}
\label{sec:payoff}

The range $\lambda' > 1/2 + o(1)$ lies outside the proven
$\vartheta < 1$ ladder regime; the entries below
are payoff-curve data for the moderate-deviation target
$\mathrm{MD}(\lambda, \eta)$ --- the open problem of upgrading the
fixed-$\lambda$ divergent-cutoff spectral-tail statement, available for every
$\lambda \in (1/2, 2/3)$ only above $V_0 = T^{2\lambda - 1 + \varepsilon}$, to an
all-$V$ count ladder $n(V) \le \Cled N V^{-4}$ valid down to $V \ge T^\eta$ ---
not certificate claims, and column generation ran at reduced budget, so the
positive values are heuristic and upper-anchored.

\begin{table}[h]
\centering
\begin{tabular}{@{}lllll@{}}
\toprule
$\lambda'$ & $\varepsilon$ & $P_{\mathrm{base}}$ & $P_{\mathrm{full}}$ &
$\delta_0$ \\
\midrule
$0.55$ & $0.05$ & $0.8053310433508489$ & $0.8053310433508484$ & $0$ \\
$0.55$ & $0.02$ & $0.8251787278264935$ & $0.8251787278264899$ & $0$ \\
$0.60$ & $0.05$ & $0.8058514020011891$ & $0.8062187975519491$ &
$+3.67 \times 10^{-4}$ \\
$0.60$ & $0.02$ & $0.8256842190868753$ & $0.8259625547887860$ &
$+2.78 \times 10^{-4}$ \\
$0.65$ & $0.05$ & $0.8059503048717801$ & $0.8086104956189072$ &
$+2.66 \times 10^{-3}$ \\
$0.65$ & $0.02$ & $0.8263688470852664$ & $0.8278795159840451$ &
$+1.51 \times 10^{-3}$ \\
\bottomrule
\end{tabular}
\end{table}

The cubic row's marginal value turns positive between $\lambda' = 0.55$ and
$0.60$ --- precisely where the flat-window moment margin
$\Marg(\lambda) = 2 m_2 - m_3$ changes sign (Corollary~\ref{cor:anchors}:
$\lambda_0 = 0.610511$) --- and grows toward the Rudnick--Sarnak endpoint
$\lambda' = 2/3$. The first configuration that could pay therefore sits at
$\lambda' \gtrsim 0.6$; consuming it would require extending the proven
large-value ladder into exactly the moderate-deviation regime
$\mathrm{MD}(\lambda, \eta)$, which is left open here as a community
target. At the proven operating point $\lambda' = 1/2 + o(1)$, the value is
exactly zero, which is what this paper proves; the curve above marks the
boundary of that statement.

\section{Consequences}
\label{sec:consequences}

\begin{enumerate}
\item \textbf{\texttt{lemmaR\_tight} stands strengthened.} The two-moment
  degeneracy is not broken; it now carries an exact-zero certificate:
  two-bandwidth flat-window augmentation --- $\lambda'$-Frobenius, signed cubic,
  all-$V$ ladder, capacity fuzz, at $\lambda' = 1/2$ --- adds exactly zero to the
  two-moment optimum at $N = 64/128$ scale, for both benchmarks, including inside
  the near-CUE class (by slack). This upgrades the standing conclusion from
  ``the two-moment certificate is exhausted'' to ``exhausted, and robust under
  Rudnick--Sarnak-range cubic augmentation with capacity control.''
\item \textbf{The $V$-dependent bridge is not warranted by the cubic payoff
  alone.} What that bridge would buy for the cubic block is $\delta_0 = 0$. The
  all-$V$ spectral-tail
  statement keeps its standalone value as the first such statement for the
  Weil--Gabor matrix, and its other consumers are unaffected.
\item \textbf{What is left standing, and what remains open, along this line:}
  (i) the all-$V$ spectral-tail statement as a standalone analytic result
  (all-$V$ count ladder with explicit
  $\Cled$, $\vartheta < 1$ strict, $C^3$ taper); separately, its fixed-$\lambda$
  divergent-cutoff form,
  $V_0 = T^{2\lambda - 1 + \varepsilon}$, extends to all
  $\lambda \in (1/2, 2/3)$ via the fourth-moment Montgomery--Vaughan branch alone
  --- a divergent-cutoff
  statement, subject to Theorem~\ref{thm:garnish}'s vacuity, not an all-$V$
  ladder; (ii) partial orthonormalization off the $o(N)$-dimensional exceptional
  subspace, where a divergent cutoff suffices (unlike the cubic LP) ---
  rank-perturbation bookkeeping already in the formalized linear-algebra library;
  (iii) the near-critical observation $X^2/T = (\log T)^{2\vartheta} = $ polylog,
  shrinking the shift range needed by quartic Hardy--Littlewood-type input for
  shifted prime pairs to logarithmic --- a possible route via averaged
  Goldston--Montgomery / Montgomery--Soundararajan results; (iv)
  $\mathrm{MD}(\lambda, \eta)$ as a community target, now carrying the
  quantified payoff curve of Section~\ref{sec:payoff}.
\item \textbf{Methodological consequence.} The instrument is reusable
  independently of the answer it returned here: an explicit-configuration LP
  whose configurations are realizable by construction, clustered null
  budgets, an explicitly fixed decision rule and stopping criterion, absorption
  claims backed by an exhibited witness or not made at all, and the robustness
  checks of Section~\ref{sec:certification}. Theorem~\ref{thm:blind} shows what
  the structural half of that instrument buys: it refutes the ``$+8$ vs $0$''
  doubles-versus-pairs separation premise outright, with no LP involved.
\end{enumerate}

\section{Formalization: what is machine-checked, and what remains}
\label{sec:formalization}

The core of Sections~\ref{sec:structure} and~\ref{sec:absorption} is
formalized, compiled, and machine-checked in Lean~4 against Mathlib, with no
\texttt{sorry}, no \texttt{native\_decide}, and no numeric certificates. Three
files under \texttt{lean/Zeta23/}\allowbreak\texttt{PairCeiling/} of the
accompanying repository carry it (Lean 4.33.0-rc2; Mathlib rev
\texttt{51e6992}; the build completes successfully):

\begin{itemize}
\item \textbf{\texttt{GridParseval.lean}.} \texttt{flat\_band\_trace\_sq} ---
  Lemma~\ref{lem:flatband}, proved over any integral domain carrying a
  primitive $M$-th root of unity, for any band of $M$ consecutive harmonics; and
  \texttt{trace\_sq\_grid} --- Theorem~\ref{thm:parseval}, the literal
  $\tr \Gh^2 = \sum_k m_k^2$ collapse on the grid. Supporting: character
  orthogonality on $\Z/M$, DFT Parseval in the integer-mark pairing, the band
  reindexing, and the triangular-weight telescoping.
\item \textbf{\texttt{GridWitness.lean}.} \texttt{vacancy\_F1} --- the $N = 64$
  vacancy lattice has $F_1 = \sum m^2 = 64$ for \emph{every} vacancy position;
  \texttt{vacancy\_half\_band\_value} --- its $\lambda' = 1/2$ half-band row is
  exactly $F' = 4128/33$, again for every vacancy position, the hand-checkable
  worked value of Section~\ref{sec:decoupling} and the witness column, its
  integer core $136224$ checked by \texttt{decide +kernel} (kernel-checked, not
  \texttt{native\_decide}); and \texttt{half\_band\_alive} ---
  $F' = 4128/33 \ne 64 = \sum m^2$, which is the ``alive'' half of
  Proposition~\ref{prop:halfband}, i.e.\ the two-bandwidth decoupling itself, in
  witness form. The finite/algebraic half of Proposition~\ref{prop:halfband}'s
  display is formalized as \texttt{sum\_band\_pair\_tri} /
  \texttt{half\_band\_fold}, together with \texttt{sum\_window\_residues} ---
  that the $s$-window is a complete residue system mod $65$, so no
  residue-class telescoping can occur, unlike the $\lambda = 1$ band.
\item \textbf{\texttt{GridCorner.lean}.} \texttt{mark\_one\_count\_ge} ---
  Lemma~\ref{lem:cornerineq} at the simple-fraction level;
  \texttt{grid\_corner\_pointwise} --- Theorem~\ref{thm:corner} pointwise on the
  grid class; \texttt{grid\_corner\_law} --- the same in law form, i.e.\ the LP
  lower bound with no discretization; and \texttt{grid\_corner\_attained} ---
  exact attainment at $N = 64$, $\varepsilon = 1/32$ by a marks-$\{1, 2\}$
  column ($12$ doubles, $40$ simples), meeting both corners with equality
  simultaneously.
\end{itemize}

\paragraph{Axiom footprint.}
\texttt{\#print axioms} on each of the twelve theorems named above returns
\texttt{[propext, Classical.choice, Quot.sound]} and nothing else --- the three
axioms of Lean's classical foundation. In particular there is no
\texttt{sorryAx}, which any admitted step would introduce, and no
\texttt{Lean.ofReduceBool}, which \texttt{native\_decide} would introduce and
which would put the compiler's evaluator into the trusted base. Full record,
with the environment and the reproduction recipe:
\texttt{results/a4-no-go/}\allowbreak\texttt{formalization-status.md} in the
accompanying repository (see \emph{Provenance, data and code}).

The formalization reaches this far for a structural reason: on the
$\psi_1$-zero grid every datum is algebraic, and Theorem~\ref{thm:parseval}
turns the first step of the corner chain from an inequality into an equality
($F_1 = \sum m^2$ exactly), after which the two integrality levels finish in
finite arithmetic.

\paragraph{What is not formalized, and what each would take.}

\begin{itemize}
\item \textbf{Proposition~\ref{prop:halfband}'s general position-sensitivity
  statement} --- two grid configurations with the same mark multiset but
  different $F'$. Finite algebra at fixed $N$, but it needs a cyclotomic
  non-vanishing certificate; the precise gap is documented in
  \texttt{GridWitness.lean}'s header. The witness form is proved; the general
  form is not.
\item \textbf{Theorem~\ref{thm:corner} off the grid} (atom-only, arbitrary
  positions) --- needs the position-space form of Lemma~\ref{lem:diagpos},
  i.e.\ $F_1 \ge \sum m^2$ with the kernel $K = D^2$ psd. This is the one step
  the grid case gets for free.
\item \textbf{Proposition~\ref{prop:affine} and Theorem~\ref{thm:blind}} ---
  polynomial identities; \texttt{ring}-adjacent, and the cheapest remaining
  items.
\item \textbf{Theorem~\ref{thm:closedforms}} --- piecewise-polynomial
  integration against Mathlib's integral.
\item \textbf{Theorem~\ref{thm:backstop}} --- a finite Hermitian matrix, inertia
  subadditivity, one scalar inequality; formalization-friendly, and the natural
  next unit from the pair channel.
\item \textbf{Theorem~\ref{thm:capacity} and Theorem~\ref{thm:garnish}} --- one
  page of real analysis (Tonelli plus a one-variable optimization) and a
  bookkeeping argument respectively.
\item \textbf{The pair-channel capacity results of
  Section~\ref{sec:pairchannel}} --- Python-certified, with interval arithmetic.
  Formalizing the certified curves is a larger undertaking than the rest
  combined.
\item \textbf{The full LP witness with its dual certificate} --- the
  exact-rational route is available: the witness law is three columns with
  rational weights, integer site subsets and marks in $\{1, 2\}$, every row value
  rational ($F_1 = 64, 96, 128$; $F'$ and $C'$ rational combinations of the
  triangular weights, e.g.\ $4128/33$), so feasibility-plus-optimality is an
  exact-rational check in the LawN256/CeilingLaw256 architecture with one added
  row family. The near-CUE witnesses extend the same format --- the pinning rows
  are rational box constraints on the stored form factors, and both witnesses
  already pass exact \texttt{Fraction} re-verification against the true
  $\tau_2 = 1$ box.
\end{itemize}

\appendix

\section{The witness data}
\label{app:witness}

\paragraph{The primary witness}
(\texttt{results/a4-m2-gate/}\allowbreak\texttt{witness\_N64.json}; cell: $N = 64$, matched budgets
$B_1, B_2, B_3 = 84.71214548308222$, $135.09544203420174$,
$304.5549106659001$, $\varepsilon = 0.05$, fuzz $=$ none). All three columns live
on the $65$-site $\psi_1$-zero grid, sites $\theta_k = k\,(64/65)$, marks
$\{1, 2\}$, no pairs. Site subsets, summarized (full integer site lists are
recorded per column in the JSON, keys \texttt{law[].atoms} /
\texttt{grid\_sites}; the canonical description is the clean integer grid --- the
stored \texttt{clean\_checks} record the maximum row-value deviation between the
clean-grid and optimizer values of $F_1/F'/C'$, at most $4.1 \times 10^{-5}$, per
column $3.4 \times 10^{-6} / 1.6 \times 10^{-5} / 4.1 \times 10^{-5}$):

\begin{itemize}
\item Column 1 (weight $0.5299906673474316$): the vacancy lattice --- $64$
  simples on $64$ of the $65$ sites (one vacancy). $\Nd = 64$; $F_1 = 64$
  exactly; $F' = 4128/33$ exactly (hand derivation: $4096/33 + 32/33$,
  Section~\ref{sec:decoupling}); $C' = 245.4508724$.
\item Column 2 (weight $0.16040139164388625$): $16$ doubles $+\ 32$ simples on
  an explicit $48$-site subset. $\Nd = 48$; $F_1 = 96$ exactly;
  $F' = 153.7604346$; $C' = 411.0865130$.
\item Column 3 (weight $0.30960794100868216$): $32$ doubles on an explicit
  $32$-site subset. $\Nd = 32$; $F_1 = 128$ exactly; $F' = 135.1959163$;
  $C' = 301.3541285$ (clean-grid values; the stored law row has
  $C' = 301.3541693$).
\end{itemize}

Aggregates (computed from the stored law rows, which saturate the budget edges
exactly): $E[F_1] = 88.94775275723633 = B_1(1 + \varepsilon)$;
$E[F'] = 132.81813304702743 \in [128.34, 141.85]$;
$E[C'] = 289.32716513260516 = B_3(1 - \varepsilon)$; $E[\Nd] = 51.52612362142$,
$E[\Nd]/N = 0.8050956815846875$. Active rows: \texttt{F1\_hi},
\texttt{Cp\_lo} only. Exact \texttt{Fraction} re-verification: all $10$ row checks
true, objective deviation $5 \times 10^{-13}$.

\paragraph{Near-CUE witnesses}
(\texttt{results/a4-m2-gate/}\allowbreak\texttt{followup-near-cue.json}, keys
\texttt{witness\_N64}, \texttt{witness\_N128}). $N = 64$: $64$ columns, marks
$\{1, 2\}$, no pairs, weights $6.2 \times 10^{-5}$--$0.0753$, $\Nd$ per column
$50$--$60$ ($4$--$14$ doubles on CUE-like positions); $E|c_j|^2 = j + 1$ at every
$j = 1, \ldots, 63$ (all upper pinning rows active at $+\tau_2$);
$E[\Nd]/N = 53.32693798297409/64 = 0.8332334059839699$; only pinning rows (and
mass) active. $N = 128$: $128$ columns, same class, upper pinning edge at
$j + 1 - 10^{-5}$ (interior shrink; \texttt{Fraction} re-verification against the
true $\tau_2 = 1$ box passes, margin $+1.0 \times 10^{-5}$);
$E[\Nd]/N = 107.19000239208492/128 = 0.8374219$ (the witness was solved at the
interior-shrunk box; its objective sits $4.0 \times 10^{-8}$ above the
Section~\ref{sec:t2result} LP optimum $0.8374218534574983$);
$E[n_{\mathrm{simple}}]/N = 0.6748$.

The three JSON files (\texttt{witness\_N64.json},
\texttt{followup-near-cue.json}, \texttt{followup-p1.json}) are in the repository
cited in \emph{Provenance, data and code}, under
\texttt{results/a4-m2-gate/}: each contains the full column data, law weights,
per-row values, active-row lists, and the rational-verification records.

\section{Null-model Monte Carlo tables}
\label{app:mc}

Closed-form anchors (Theorem~\ref{thm:closedforms}): $m_2(1) = 4/3$,
$m_2(1/2) = 13/6 = 2.166667$, $m_3(1/2) = 5$.

\paragraph{Primary sampler}
(CUE eigenangles via QR-with-phase-fix of complex Ginibre, unfolded to
circumference $n$; $R$ adaptive to $\mathrm{SE}(m_3) \le 0.5\%$; batch-means SE
with jackknife check):

\begin{table}[h]
\centering
\begin{tabular}{@{}lllll@{}}
\toprule
$n$ & $R$ & $m_2(1)$ & $m_2(1/2)$ & $m_3(1/2)$ \\
\midrule
$64$ & $4000$ & $1.32363(74)$ & $2.11087(63)$ & $4.75867(362)$ \\
$128$ & $2000$ & $1.32867(77)$ & $2.13845(71)$ & $4.87769(411)$ \\
$256$ & $600$ & $1.33115(66)$ & $2.15231(63)$ & $4.93758(373)$ \\
$512$ & $300$ & $1.33199(86)$ & $2.15958(67)$ & $4.96843(413)$ \\
$1024$ & $120$ & $1.33364(110)$ & $2.16362(83)$ & $4.98715(476)$ \\
$2048$ & $48$ & $1.33400(94)$ & $2.16505(83)$ & $4.99323(504)$ \\
\addlinespace
$1/n$ fit $\to \infty$ & & $1.33394(63)$ & $2.16682(52)$ & $5.00056(311)$ \\
closed form & & $1.33333$ & $2.16667$ & $5$ \\
\bottomrule
\end{tabular}
\end{table}

All three extrapolations are consistent with the closed-form anchors.

\paragraph{Independent sampler}
(independent QR sampler, fresh seeds), side by side at the two shared sizes ---
agreement within ${\sim}1$ combined $\sigma$ per entry:

\begin{table}[h]
\centering
\begin{tabular}{@{}llll@{}}
\toprule
$n$ & $m_2(1)$ indep. / primary & $m_2(1/2)$ indep. / primary &
$m_3(1/2)$ indep. / primary \\
\midrule
$64$ & $1.32317(105) / 1.32363(74)$ & $2.11120(84) / 2.11087(63)$ &
$4.76099(491) / 4.75867(362)$ \\
$256$ & $1.33045(132) / 1.33115(66)$ & $2.15305(103) / 2.15231(63)$ &
$4.94059(616) / 4.93758(373)$ \\
\bottomrule
\end{tabular}
\end{table}

A third independent sampler (own seeds, alias-free
frequency-matrix assembly) gave $m_2(1/2) = 2.11063(135)$ at $n = 64$ and
$2.15273(100)$ at $n = 256$, $m_3(1/2) = 4.7589(79)$ and $4.9408(59)$ --- each
within ${\sim}1\sigma$ of both prior samplers --- with two-point $1/n$
extrapolation $m_2 \to 2.1668$, $m_3 \to 5.0014$. A third route for
$m_3(1/2) = 5$ (real-space 2D determinantal quadrature) gave
$4.9896 \to 4.9932$ under truncation refinement; a further independent
quadrature gave $4.9924 \to 4.9960$.

\paragraph{Ladder and occupancy side data}
(matched $n = 64$ unless noted): null ladder floor
$\Cled^{\min} = 3.8188$ ($3.9119$ at $n = 128$);
$E[n(V)]/n$ on $V = [2, 3, 4, \ldots] = [0.23868, 0.01364, 0, \ldots]$;
$\max |\mathrm{eig}|$ over $4000$ draws $= 3.977$. Cross-term (occupancy $\ge 2$)
shares: Frobenius $0.52626$, cubic $0.78986$, against the analytic $54\%$ ($7/6$
of $13/6$) and $80\%$ ($4$ of $5$; of which the strictly two-point share is $3.5$
and the three-point term $T_3$ the remaining $0.5$).

\section*{Provenance, data and code}
\addcontentsline{toc}{section}{Provenance, data and code}

\paragraph{Method.}
The specification of the decision computation --- objective, adversary class,
constraint rows, decision rule and stopping criterion --- was fixed in writing
before the computation was run. That specification, the run record and the
verification outputs are among the artifacts listed below.

\paragraph{Data and code availability.}
Every number in this paper comes from one of three places: the decision record,
the verification suite written for this paper, or the Lean development. All
three are version-controlled and are available in the accompanying
repository \url{https://github.com/x67ai/riemann-rh-program}, at the tag
\texttt{paper-2026-08-28}, which is the state this paper was prepared against.
Paths below are relative to \texttt{rh-program/} in that repository. One
exception: the working PDF corpus of fetched literature is not redistributed,
for copyright reasons, and nothing in this paper depends on it --- every
external claim is cited to its published source. The components are:

\begin{itemize}
\item \emph{The formalization record}
  (\texttt{results/a4-no-go/formalization-status.md}) behind
  Section~\ref{sec:formalization}: the environment (Lean 4.33.0-rc2, Mathlib rev
  \texttt{51e6992}), the build result, the theorem-by-theorem map to this
  paper's numbering, the \texttt{\#print axioms} output for all twelve
  machine-checked theorems, and the reproduction recipe.
\item \emph{The specification of the corrected row system}
  (\texttt{results/adjudication-A4.json}): the garnish and squeeze
  computations, and the row system in its corrected form, against which the
  computation was run.
\item \emph{The specification of the decision computation}
  (\texttt{results/a4-m2-gate/SPEC.md}): kernels and units, null budgets and
  their closed forms with derivations, the adversary class and its
  realizability argument, the row system and capacity constant, the decision
  rules, and the flags.
\item \emph{The run report} (\texttt{results/a4-m2-gate/RUN-REPORT.md}): the
  outcome, the grid-decoupling mechanism, the sanity checks, every run and
  number produced, the witness, and every deviation from the specification.
\item \emph{The independent check} (\texttt{results/a4-m2-gate/AUDIT.md}): a
  Monte-Carlo check of the null budgets, an independent re-derivation of the
  budget identities and the block law, recomputation of the witness, LP
  re-solves over all $4003$ columns, and eight robustness checks of the
  conclusion along independent axes.
\item \emph{The follow-up report}
  (\texttt{results/a4-m2-gate/FOLLOWUP-REPORT.md}): the near-CUE runs at
  $\tau_2 = 1$ for both $N$, the $p_1$ objective, and the riders attached to the
  use of~\cite{BGSTB24}.
\item \emph{The raw data} (\texttt{results/a4-m2-gate/*.json},
\texttt{results/a4-m2-gate/}\allowbreak\texttt{runs/*.json}): the decision-LP
result, the witness law at $N = 64$, the
  near-CUE and $p_1$ follow-up outputs, and the full scan records at $N = 64$ and
  $N = 128$ (including the tight-$\varepsilon$, near-CUE, window and
  $\lambda'$-exploration runs, and the $\lambda' = 1/2$ budget derivations), all
  as JSON.
\item \emph{The proof note} (\texttt{results/a4-no-go/theorems.md}): the six
  proofs incorporated as Sections~\ref{sec:structure}, \ref{sec:atomonly}
  and~\ref{sec:capacity}, with their verification scripts and JSON outputs
  (\texttt{results/a4-no-go/verify/}).
\item \emph{The pair-channel note}
  (\texttt{results/a4-no-go/}\allowbreak\texttt{pair-}\allowbreak\texttt{channel.md}):
  the pair-interference closure
  (Sections~\ref{sec:fractional}--\ref{sec:pairchannel}), including the proofs
  of Theorems~\ref{thm:singlepair} and~\ref{thm:multipair}, for which only the
  proof architecture is given here, with its independent verification suite
  \texttt{results/a4-no-go/verify/}\allowbreak\texttt{pairchan\_*.py}
  and \texttt{pairchan\_}\allowbreak\texttt{verify\_}\allowbreak\texttt{out.json}.
\item \emph{The data digest} (\texttt{results/a4-no-go/data-tables.md}): the
  digest from which every number in this paper is cited, each row carrying its
  on-disk source (file plus JSON key or report section);
  its Section~10 records two numerical cautions observed here (the $940$-record
  maximum is the $N = 128$ value $8.94 \times 10^{-13}$; the near-CUE level
  $0.833377 / 0.8333622 / 0.8332334$ drift across dictionary snapshots is
  common-mode, the equality $\delta_0' = 0$ being snapshot-independent).
\item \emph{The proposal record} (\texttt{directions/A4-lindelof-lock.md}): the
  original proposal and the specification of the computation reported here.
\end{itemize}

\begin{thebibliography}{99}

\bibitem[AF26]{AF26} L. Alp\"oge and R. Furman, \emph{More than two thirds of
the zeta zeros are simple and on the critical line}, arXiv:2608.13637 (v1, 13~Aug~2026;
v2, 19~Aug~2026 --- the arXiv record carries these two versions and no others),
with its sorry-free Lean~4 formalization
(Theorems A--E; \texttt{lemmaR\_tight} in
\texttt{Zeta23/ZeroSide/TightMult.lean}; the PairCeiling library,
LawN256/CeilingLaw256; Theorem~5.8/Remark~5.10 prime-side moments).
Two separately circulated revisions, neither of them an arXiv version, differ in
their sectioning, and this paper cites them separately:
(i) the 35-page version, \emph{More than two thirds of the zeros of the Riemann
zeta function lie on the critical line}, whose
\S7.5 ``Limits of the method'' carries the item (e) quoted in
Section~1.2; and
(ii) the 17-page version, called v5 here, \emph{\ldots are simple and on the
critical line}, in which the same material appears,
in a shorter form, as \S7.2(e), \S7.2 there being titled ``Sharpness and limits
of the method''. Where this paper writes \S7.5(e) it means (i); where it writes
\S7.2(e) it means (ii). Both readings are given in Section~1.2.

\bibitem[BGSTB24]{BGSTB24} S.~A.~C. Baluyot, D.~A. Goldston, A.~I. Suriajaya,
C.~L. Turnage-Butterbaugh, \emph{An unconditional Montgomery theorem for pair
correlation of zeros of the Riemann zeta function}, arXiv:2306.04799
(v1, 7~Jun~2023); Acta Arith. \textbf{214} (2024), 357--376;
doi:10.4064/aa230612-20-3. Cited for: Theorem~1 (unconditional pointwise
$F(\alpha)$ on the closed band $0 \le \alpha \le 1$, depth-weighted,
multiplicity-counted), the Lemma~5 kernel class, and the Remark after
Theorem~2 (rider~1). See also the same authors' sequel on the horizontal
distribution of zeros, \emph{Pair correlation of zeros of the Riemann zeta
function I: proportions of simple zeros and critical zeros}, arXiv:2501.14545
(v1, 24~Jan~2025; v2, 21~Nov~2025; unpublished as of August~2026), which works
under a narrow-vertical-box hypothesis.

\bibitem[LR20]{LR20} J.~C. Lagarias and B. Rodgers, \emph{Higher correlations
and the Alternative Hypothesis}, Q.~J. Math. \textbf{71} (2020), no.~1,
257--280; doi:10.1093/qmathj/haz043; arXiv:1905.12123. Cited in
Section~1.5: they prove that the currently known band-limited $n$-level
correlation data for zeta zeros --- Montgomery ($n = 2$), Hejhal ($n = 3$),
Rudnick--Sarnak ($n \ge 4$) --- does not rule out the Alternative
Hypothesis, by exhibiting a point process on $(1/2)\Z$ matching the sine process
against all such test functions. Tao obtained the same conclusion independently,
by different methods (cited there).

\bibitem[LR21]{LR21} J.~C. Lagarias and B. Rodgers, \emph{Band-limited mimicry
of point processes by point processes supported on a lattice}, Ann. Appl.
Probab. \textbf{31} (2021), no.~1, 351--376; doi:10.1214/20-AAP1592;
arXiv:1907.03391. Their Theorem~1.4 and Proposition~3.4 show that above the
Nyquist bandwidth the mimicking correlations are unique, while at
critical sampling they are not --- the same degeneracy this paper's
Theorem~\ref{thm:parseval} exploits, in the same place; the consequences for the
scope of the absorption are discussed at the point of citation in
Section~\ref{sec:decoupling}.

\bibitem[Mon73]{Mon73} H.~L. Montgomery, \emph{The pair correlation of zeros of
the zeta function}, in Analytic Number Theory (St.~Louis, 1972), Proc. Sympos.
Pure Math. 24, Amer. Math. Soc., 1973, 181--193. The source of the extremal
values $5/6$ and $2/3$ whose bandwidth-one analogs appear in
Theorem~\ref{thm:corner} and Corollary~\ref{cor:matched}.

\bibitem[Hej94]{Hej94} D.~A. Hejhal, \emph{On the triple correlation of zeros of
the zeta function}, Internat. Math. Res. Notices \textbf{1994}, no.~7,
293--302; doi:10.1155/S1073792894000334. The $n = 3$ correlation input named in
Section~1.2.

\bibitem[RS96]{RS96} Z. Rudnick and P. Sarnak, \emph{Zeros of principal
$L$-functions and random matrix theory}, Duke Math. J. \textbf{81} (1996),
no.~2, 269--322; doi:10.1215/S0012-7094-96-08115-6. The $n$-level range that
names this
paper's operating window. Scope, relevant to the cubic budget center
(Section~\ref{sec:flags}): their unconditional $n$-level theorem is for smooth
test functions; taking a characteristic function requires RH.
Section~\ref{sec:sine}'s determinantal $\rho_3$ and the $c_3(\lambda')$
identification invoke the smooth statement.

\bibitem[CGG98]{CGG98} J.~B. Conrey, A. Ghosh and S.~M. Gonek, \emph{Simple
zeros of the Riemann zeta-function}, Proc. London Math. Soc. (3) \textbf{76}
(1998), 497--522; doi:10.1112/S0024611598000306. The origin of the second
integrality level that
Lemma~\ref{lem:cornerineq} and Theorem~\ref{thm:corner} consume, made
unconditional in Alp\"oge and Furman's \S7.2(c).

\bibitem[BHB13]{BHB13} H.~M. Bui and D.~R. Heath-Brown, \emph{On simple zeros of
the Riemann zeta-function}, Bull. London Math. Soc. \textbf{45} (2013),
no.~5, 953--961; doi:10.1112/blms/bdt026; arXiv:1302.5018. The conditional
simple-zero input quoted in Section~1.2; their unconditional-on-RH
distinct-zeros constant, as published, is $0.84665$.

\bibitem[CGdL20]{CGdL20} A. Chirre, F. Gon\c{c}alves and D. de Laat, \emph{Pair
correlation estimates for the zeros of the zeta function via semidefinite
programming}, Adv. Math. \textbf{361} (2020), 106926;
doi:10.1016/j.aim.2019.106926; arXiv:1810.08843v2. The
semidefinite-programming line for pair-correlation extremal problems, of which
this paper's LP is the bandwidth-one marked-configuration analog. Two specific
contacts. (i) Their inequality~(7),
$2 N_s(T) \le \sum_\rho (m_\rho - 2)(m_\rho - 3)/m_\rho
 = N^{*}(T) - 5N(T) + 6\Nd(T)$, credited there to an argument of Ghosh, is the
published instance of the same per-multiplicity integrality device
Lemma~\ref{lem:cornerineq} uses, at the reciprocal-weighted third level
$(m-2)(m-3) \ge 0$; Lemma~\ref{lem:cornerineq}(a) is the unweighted
second level $(m-1)(m-2) \ge 0$. (Their semidefinite programming is applied to the
auxiliary-function class of their Section~3, not to~(7); (7) is elementary integer
arithmetic.) (ii) Their Corollary~3 is the corresponding published record for
distinct zeros: $\Nd(T) \ge (0.8477 + o(1))N(T)$ on RH and
$(0.8486 + o(1))N(T)$ on GRH.

\bibitem[GdLL25]{GdLL25} F. Gon\c{c}alves, D. de Laat and N. Leijenhorst,
\emph{Multiplicity of nontrivial zeros of primitive $L$-functions via higher-level
correlations}, Math. Comp. \textbf{94} (2025), no.~354, 2041--2058;
doi:10.1090/mcom/4005; arXiv:2303.01095. The closest published third-party work to
this paper: the same higher-level-correlation-versus-multiplicity question,
attacked
by semidefinite programming rather than by a marked-configuration LP. Cited in
Sections~1.2 and~3.2. Their equation~(5),
$N_n(T) = \sum_k S(n,k) M_k(T)$ with $S(n,k)$ the Stirling numbers of the second
kind, is the change of basis whose $n = 3$ case is the atom half of
Theorem~\ref{thm:blind} here.

\bibitem[FGL14]{FGL14} D.~W. Farmer, S.~M. Gonek and Y. Lee, \emph{Pair correlation
of the zeros of the derivative of the Riemann $\xi$-function}, J. London Math.
Soc.~(2) \textbf{90} (2014), no.~1, 241--269; doi:10.1112/jlms/jdu026. Cited
only in
Corollary~\ref{cor:matched}, to disambiguate a numerical coincidence: their
published
RH-conditional lower bound for the proportion of distinct zeros is
$\Nd(T) > 0.8051\,N(T)$ (their Corollary~1.4). The two-author 2008 preprint
arXiv:0803.0425 is an earlier version whose Corollary~1.4 gives $0.6544$; the
constant
$0.8051$ is the published version's.

\bibitem[KLS07]{KLS07} T. Kuna, J.~L. Lebowitz and E.~R. Speer,
\emph{Realizability of point processes}, J. Stat. Phys. \textbf{129} (2007),
no.~3, 417--439; doi:10.1007/s10955-007-9393-y;
arXiv:math-ph/0612075. The truncated-moment/realizability problem that
Section~\ref{sec:adversary}'s adversary-realizability argument re-derives in the
finite-circle setting.

\bibitem[KLS11]{KLS11} T. Kuna, J.~L. Lebowitz and E.~R. Speer,
\emph{Necessary and sufficient conditions for realizability of point processes},
Ann. Appl. Probab. \textbf{21} (2011), no.~4, 1253--1281; doi:10.1214/10-AAP703;
arXiv:0910.1710. The sequel, which settles the general finite-order
truncated-moment
realizability question; Section~\ref{sec:adversary}'s finite-circle argument is a
special case in their setting rather than a new problem.

\bibitem[CdLS22]{CdLS22} H. Cohn, D. de Laat and A. Salmon, \emph{Three-point
bounds for sphere packing}, arXiv:2206.15373 (v1, 30~Jun~2022; unpublished as
of August~2026). The comparandum: in the
sphere-packing analog the third-order augmentation \emph{does} pay, improving
on the two-point Cohn--Elkies bound in several dimensions; here it is worth
exactly nothing.

\bibitem[KS66]{KS66} S. Karlin and W.~J. Studden, \emph{Tchebycheff Systems:
With Applications in Analysis and Statistics}, Pure and Applied Mathematics,
Vol.~XV, Interscience Publishers (John Wiley \& Sons), New York, 1966. The
classical Chebyshev--Markov--Stieltjes/Christoffel background for the fact that
odd moments do not improve one-sided bounds --- the general form of the cubic
blindness of Theorem~\ref{thm:blind}.

\end{thebibliography}

\end{document}
